\documentclass[11pt,a4paper]{article}
\usepackage[T1]{fontenc}
\usepackage[utf8]{inputenc}
\usepackage{lmodern,microtype}
\usepackage[a4paper,margin=27mm,headheight=15pt]{geometry}
\usepackage{amsmath,amssymb,amsthm,mathtools}
\usepackage{booktabs,longtable,array,enumitem}
\newcolumntype{P}[1]{>{\raggedright\arraybackslash}p{#1}}
\usepackage[dvipsnames]{xcolor}
\usepackage{fancyhdr}
\usepackage[unicode,breaklinks,colorlinks=true,linkcolor=MidnightBlue,citecolor=MidnightBlue,urlcolor=MidnightBlue]{hyperref}
\usepackage[nameinlink,noabbrev]{cleveref}
\usepackage{listings}
\hypersetup{pdftitle={A Cantor Space of Surreal Subfields},pdfauthor={AI-assisted research article prepared for Vladimir Reshetnikov},pdfsubject={Finite supports, complete analytic embeddability, and generic behavior of countable surreal subfields}}
\pagestyle{fancy}
\fancyhf{}
\fancyhead[L]{\small\textsc{A Cantor Space of Surreal Subfields}}
\fancyhead[R]{\small Research draft}
\fancyfoot[C]{\thepage}
\renewcommand{\headrulewidth}{0.3pt}
\setlength{\parskip}{4pt}
\setlength{\parindent}{1.2em}
\setlength{\emergencystretch}{2em}
\setlist{itemsep=3pt,topsep=5pt}
\numberwithin{equation}{section}
\newtheorem{theorem}{Theorem}[section]
\newtheorem{proposition}[theorem]{Proposition}
\newtheorem{lemma}[theorem]{Lemma}
\newtheorem{corollary}[theorem]{Corollary}
\theoremstyle{definition}
\newtheorem{definition}[theorem]{Definition}
\newtheorem{example}[theorem]{Example}
\newtheorem{question}[theorem]{Question}
\theoremstyle{remark}
\newtheorem{remark}[theorem]{Remark}
\newcommand{\N}{\mathbb N}
\newcommand{\Q}{\mathbb Q}
\newcommand{\R}{\mathbb R}
\newcommand{\Z}{\mathbb Z}
\newcommand{\No}{\mathbf{No}}
\newcommand{\Ord}{\mathrm{Ord}}
\newcommand{\rc}{\operatorname{rc}}
\newcommand{\cl}{\operatorname{cl}}
\newcommand{\trdeg}{\operatorname{trdeg}}
\newcommand{\supp}{\operatorname{supp}}
\newcommand{\rank}{\operatorname{rk}}
\newcommand{\emb}{\hookrightarrow}
\newcommand{\Cstar}{\mathcal C^{\star}}
\newcommand{\Cone}{\mathcal C^{(1)}}
\newcommand{\Kstar}{K^{\star}}
\newcommand{\Kone}{K^{(1)}}
\newcommand{\Fone}{F^{(1)}}
\newcommand{\Sig}{\boldsymbol\Sigma^1_1}
\newcommand{\PiC}{\boldsymbol\Pi^1_1}
\newcommand{\LO}{\mathrm{LO}}
\newcommand{\CLO}{\mathrm{CLO}}
\newcommand{\RCF}{\mathrm{RCF}}
\newcommand{\KB}{\mathrm{KB}}
\newcommand{\IF}{\mathrm{IF}}
\newcommand{\WF}{\mathrm{WF}}
\newcommand{\Aut}{\operatorname{Aut}}
\newcommand{\ds}{\displaystyle}
\newcommand{\e}[1]{\omega^{#1}}
\newcommand{\basefield}{k}
\lstset{basicstyle=\ttfamily\small,breaklines=true,columns=fullflexible,keepspaces=true,frame=single,rulecolor=\color{black!20},showstringspaces=false}
% [write] Notes added when the manuscript joined ProveIt's surreal report
% collection (batch 90, 4 October 2026). Everything else is the delivered text,
% with every label prefixed csf: and cleveref type hints on the labels of
% lemmas, corollaries, propositions and questions (which share the theorem
% counter, so that cleveref otherwise names them all "theorem").
\crefname{question}{question}{questions}
\Crefname{question}{Question}{Questions}
\newcommand{\repo}[1]{\nolinkurl{#1}}
\newcommand{\lbl}[1]{\nolinkurl{#1}}
\newenvironment{writenote}[1][]{\par\smallskip\begingroup\small\raggedright
 \setlength{\parindent}{1.25em}\noindent\textbf{[write]}\if\relax\detokenize{#1}\relax\else~\textbf{#1.}\fi\ }%
 {\par\endgroup\smallskip}

\begin{document}
\hypersetup{pageanchor=false}
\begin{titlepage}
\vspace*{12mm}
{\color{MidnightBlue}\large\sffamily DESCRIPTIVE SET THEORY $\boldsymbol\cdot$ SURREAL FIELDS}\par
\vspace{12mm}
{\fontsize{27}{33}\selectfont\bfseries A Cantor Space\par of Surreal Subfields}\par
\vspace{7mm}
{\Large Finite supports, universal embeddings,\par and two generic regimes}\par
\vspace{10mm}
\noindent\rule{\textwidth}{0.6pt}
\vspace{5mm}
\noindent Research article prepared for Vladimir Reshetnikov\\
AI-assisted mathematical draft\hfill 3 October 2026
\vspace{11mm}
\begin{abstract}
We construct an explicit countable real-closed subfield $\Kstar$ of the surreal
numbers and a closed Cantor family of elementary subfields of $\Kstar$.
Every ambient element has a unique finite support in a fixed transcendence
basis. Consequently, varying a subfield is topologically as simple as choosing
a subset of a countable set. Nevertheless, field embeddability on this compact
family is a complete analytic quasiorder. A concrete coding uses
$\omega^{b\omega^q}$, with $q\in\Q$ and
$b\in\{1,\sqrt2,\sqrt3,\sqrt5,\ldots\}$; the ordered additive groups of real
quadratic fields intrinsically record countably many colors.

A rank-one subfamily admits an exact reduction to countable linear orders.
Within it, a fixed field of reverse-$\omega$ value-group rank has an
analytic-complete embedding locus, while every source of finite transcendence
degree has an open locus. Almost every rank-one parameter, both in category
and in coin-flip measure, gives the same isomorphism type. In the larger
family, by contrast, the locus admitting the whole ambient field is also analytic-complete
and dense, but meager and null. General category and measure zero--one laws explain
the two regimes.

The unrestricted complexity of countable real-closed-field embeddability is
already known. The contribution developed here is the explicit common
surreal ambient field, compact finite-support localization, fixed-source
threshold, and generic comparison. Full proofs of these refinements are
given; historical priority is not certified, and no Lean verification is
claimed.
\end{abstract}
\vfill
\noindent\small\textbf{Scope.} Ordinary ordered fields, not exponential or differential
fields. All descriptive-set-theoretic spaces are sets. No measurable
structure on the proper class $\No$ is used.

\smallskip
{\footnotesize\raggedright\noindent\textbf{[write]} Filed in ProveIt's surreal report collection as \repo{foundations-and-computation/cantor-families-of-surreal-subfields} (batch~90, manuscript~04; label prefix \texttt{csf:}). Provenance, formal status, the relation to neighbouring reports and a notation table are in Sections~\ref{csf:sec:provenance} and~\ref{csf:sec:notation}.\par}
\end{titlepage}
\hypersetup{pageanchor=true}

\tableofcontents
\newpage

\section{Research intersection and the precise contribution}
\label{csf:sec:intersection}

\subsection{Why this topic connects the two research programs}

Elliot Glazer's \emph{A Topological Tennenbaum Theorem} studies the tension
between algebraic structure, definability, and Polish topology in models of
arithmetic \cite{glazer}. His work on FrontierMath concerns the evaluation of
advanced mathematical reasoning \cite{frontier}. These provide two relevant
interests: descriptive obstructions to apparently simple mathematical
presentations, and the difference between genuine structural difficulty and
what a sampling procedure happens to reveal.

ProveIt has an explicit surreal-field development, real-closedness results,
Hahn-series infrastructure, and research on self-embeddings and their
valuation invariants \cite{proveit,sse}. The present article moves from
individual self-embeddings of $\No$ to a \emph{space of countable elementary
subfields of one field inside $\No$}. This connects the valuation-based
constructions with descriptive classification, without putting a Polish
space structure on all surreal numbers.

There is an important change of level. Our Cantor topology is a topology on
\emph{parameters for fields}; it is not the order topology on the domain of
one field. Each field in the family is countable. Thus none of the results
below contradicts or purports to settle Glazer's questions concerning
uncountable topological arithmetic.

\subsection{What was already known}

Calderoni, Marker, Motto Ros, and Shani proved that embeddability of countable
real closed fields is a complete analytic quasiorder, and that
bi-embeddability is a complete analytic equivalence relation
\cite[Corollary 7.3]{cmms}. Their route uses colored orders, ordered divisible
abelian groups, and monomial fields. The isomorphism complexity of
real closed fields also lies within the earlier o-minimal analysis of Rast
and Sahota \cite{rast}. We use these as precedents, not as new discoveries.

The external descriptive completeness input for our strongest theorem is
Louveau's theorem that embeddability of countably colored countable linear
orders is a complete analytic quasiorder, as recorded in
\cite[Theorem 7.1]{cmms} and \cite[Theorem 3.2]{marcone}.
Our proof reconstructs the algebraic coding explicitly and then localizes it
inside one fixed countable surreal field and a closed coordinate family.

The manuscript answers the following \emph{local refinement questions}.
Can the complexity be confined to actual subfields of one explicitly given
countable surreal field? Can this happen in a compact family with
finite-coordinate membership tests? Can a single fixed source already give
a non-Borel existence problem? How do transcendence degree and generic
sampling affect that problem? These are the questions resolved here. We do
not assert that they were previously published as open problems.

\subsection{Main results at a glance}

Let $k$ be the real algebraic numbers. Enumerate the positive rational primes
as $p_0,p_1,\ldots$, and put
\begin{equation}
 B=\{1\}\cup\{\sqrt{p_n}:n\in\N\},\qquad I=\Q\times B,
 \qquad x_{q,b}=\omega^{b\omega^q}.
 \label{csf:eq:ambient-generators}
\end{equation}
The powers of $\omega$ in this article are Conway's omega-map. In particular,
\eqref{csf:eq:ambient-generators} makes no assumption about closure under
Gonshor's exponential.

For $A\subseteq I$, let
\begin{equation}
 K_A=\rc_{\No}\bigl(k(x_{q,b}:(q,b)\in A)\bigr),
 \qquad \Kstar=K_I.
 \label{csf:eq:ambient}
\end{equation}
Here $\rc_{\No}$ means relative real closure. Its set-sized construction is
explained in \cref{csf:sec:foundations}.

\begin{theorem}[Fixed-ambient compact universality]
\label{csf:thm:main}
The field $\Kstar$ is countable and real closed. The map
\[
 \Phi:2^I\longrightarrow 2^{\Kstar},\qquad A\longmapsto K_A
\]
is a homeomorphism onto a closed Cantor family $\Cstar$. Every member of
$\Cstar$ is an elementary subfield of $\Kstar$, and all have residue field
$k$ for their natural valuations. On $\Cstar$, unital field embeddability is
a complete analytic quasiorder; bi-embeddability is a complete analytic
equivalence relation. Isomorphism is Borel complete for countable structures.
\end{theorem}

The proof separates into the finite-support theorem
(\cref{csf:thm:support,csf:thm:cantor}), the uncolored classification
(\cref{csf:thm:rank-one-classification}), and the intrinsic quadratic-color coding
(\cref{csf:thm:colored-coding,csf:thm:universality}).

For $S\subseteq\Q$, write
\[
 \Kone_S=K_{S\times\{1\}},\qquad \Kone=\Kone_{\Q},
 \qquad \Cone=\{\Kone_S:S\subseteq\Q\}.
\]
Let $D=\{-n:n\in\N\}$ with its ordinary order, so that $D$ has type
$\omega^*$. For a countable real closed field $P$, define
\[
 E^{(1)}_P=\{S\subseteq\Q:P\emb\Kone_S\},\qquad
 E^\star_P=\{A\subseteq I:P\emb K_A\}.
\]

\begin{theorem}[Fixed sources and generic behavior]
\label{csf:thm:main-source}
The following hold.
\begin{enumerate}[label=\textnormal{(\roman*)}]
\item $E^{(1)}_{\Kone_D}$ is analytic-complete; its complement is
coanalytic-complete. The hardness persists when every target has countably
infinite transcendence degree.
\item In either coordinate family, a source of finite transcendence degree
has an open embedding locus. A source of infinite transcendence degree has
a locus with empty interior. Thus a nonempty locus is open exactly in the
finite-transcendence-degree case.
\item Every nonempty $E^{(1)}_P$ is dense, comeager, and of coin-flip measure
one. In fact, a dense $G_\delta$ set of parameters gives fields isomorphic to
$\Kone$.
\item Every nonempty $E^\star_P$ is dense. Each such locus is either meager
or comeager and has measure zero or one. Both generic regimes occur:
$E^\star_{\Kone_D}$ is comeager and conull, whereas
$E^\star_{\Kstar}$ is meager and null. Both of these loci are
analytic-complete.
\end{enumerate}
\end{theorem}

All assertions are statements in ordinary ZFC. They concern actual unital
field embeddings, not only embeddings that are required in advance to
preserve chosen generators or monomials.

\subsection{Provenance and place in ProveIt}
\label{csf:sec:provenance}

\begin{writenote}[Added 4 October 2026, batch~90]
This report is a single manuscript, printed in full: manuscript~04 of batch~90 of ProveIt's incoming-reports intake, delivered as the archive \texttt{Surreal\_\allowbreak Subfield\_\allowbreak Cantor\_\allowbreak Research.zip} (544,930 bytes, nine members; inner directory \texttt{surreal\_\allowbreak subfield\_\allowbreak cantor/}, main file \texttt{article.tex}, 1843 lines, with a 27-page A4 PDF) in the arrival commit \texttt{a162e4386}, and placed in commit \texttt{12076b2e8} as \repo{Algebra/SurrealNumbers/docs/foundations-and-computation/cantor-families-of-surreal-subfields}. Its subject is real closed subfields of $\No$, so the family \texttt{surreal/} would also have fitted; it was filed under \texttt{foundations-and-computation/} because its results are descriptive set theory about subfields of $\No$, beside the complexity results of \repo{polish-models-of-omnific-arithmetic} and \repo{omnific-notations}. No report of the collection asks its questions. The author lines read ``Research article prepared for Vladimir Reshetnikov'' and ``AI-assisted mathematical draft''; the report is AI-assisted, unrefereed and not formalized, and its historical priority is not established. In the \textbf{[write]} notes, report paths are relative to \repo{Algebra/SurrealNumbers/docs/} and Lean paths to \repo{Algebra/SurrealNumbers/}, unless given in full.

\emph{Pin.} The manuscript was written against ProveIt commit \texttt{7c0f2d9f92c3d51ec85703bed1022924a4b7359b} (3~October 2026, ``New research reports''), an ancestor of the placement, quoted in \cref{csf:sec:formalization} and in \cite{proveit}. Its statements about the repository are accurate at the pin and at the write: the three Lean files of the table in \cref{csf:sec:formalization} are byte-identical at both (blobs \texttt{a0d30d424}, \texttt{632c027e9}, \texttt{e9b72519e}); \texttt{signSequenceIsRealClosed} is an \texttt{instance} (line~36 of \repo{Surreal/Foundations/SignSequenceRealClosed.lean}); \repo{Surreal/HahnSeries/RealClosed.lean} proves real closedness of set-sized Hahn fields over a real closed coefficient field with divisible exponent group; and the surreal README and the self-embedding report's README separate formal declarations from unrefereed reports (the latter README has gained cross-references since the pin, without changing that). No repository statement is contradicted, and nothing needed retraction.

\emph{Formal status.} No statement of this report is formalized. One finite algebraic input is available in Lean: \cref{csf:lem:radicals} is the instance $M=\Q$, $a_i=p_i$, square roots taken in $\R$, of \lbl{tail:lem:signs} of \repo{surreal/tail-spans-and-differential-transcendence}, which is proved as \texttt{Surreal.TailSigns.linearIndependent\_mono} (in \repo{Surreal/Algebra/MultiquadraticSigns.lean}, for any field with $2\ne0$ and roots in any extension field), with the hypothesis on the primes supplied by \texttt{Surreal.AlgebraicClosurePart.independentSquareClasses\_prime} (in \repo{Surreal/HahnSeries/AlgebraicCoefficientClosure.lean}); the specialization itself is not stated as a declaration. Nothing in the Lean tree states the double-monomial independence, the finite-support theorem, the value-group computation or any descriptive-set-theoretic result here. The formalization plan of \cref{csf:sec:formalization} is a proposal.

\emph{Review in the Hilbert's-tenth research tree.} Another session inventoried the six batch-90 archives at their arrival commit (commit \texttt{bcc1a4438}, \repo{Computability/HilbertTenthProblem/Papers/research-wip/native-stream-queue/review_polish_partx_c9bc70d8f.md}, with its JSON receipt). For this archive it read the delivered README in full (lines 1--86) and hashed the other eight members without reading them. Its routing result is that the report concerns ``Analytic-complete embeddability inside a countable surreal field and finite-support topology; descriptive complexity is not finite arithmetic universality''. It executed no program, made no correctness finding and needed no scope correction; the Hilbert's-tenth programme's 84-operation universal polynomial is unaffected.

\emph{The cited prior work.} The write checked \cite{cmms} at arXiv (version~2 of 13~January 2023; the arXiv record notes acceptance in \emph{Advances in Mathematics}). Its Section~7, ``ODAG and other o-minimal theories'', states on PDF pages 33--34 Theorem~7.1 (embeddability of countably colored linear orders is a complete analytic quasiorder, ``observed by Louveau'', citing Marcone--Rosendal, Theorem~3.2), Proposition~7.2 (a Borel reduction of colored-order embeddability to embeddability of ordered divisible abelian groups, through the finite-support reverse-lexicographic group with Archimedean component $H_{c(n)}$ at rank $n$, for pairwise non-embeddable divisible subgroups $H_n\subseteq\R$) and Corollary~7.3 (embeddability of countable real closed fields is a complete analytic quasiorder and bi-embeddability a complete analytic equivalence relation, proved on pages 34--35 through relative real closures of $k(t^g:g\in G)$ in Hahn fields); its Lemma~2.2 is Hion's lemma. The citations of \cref{csf:sec:intersection,csf:sec:colors} and of the bibliography match. \Cref{csf:thm:colored-coding,csf:thm:universality} thus compose the two steps of Proposition~7.2 and Corollary~7.3 inside one fixed field, with the explicit components $H_n=\Q(\sqrt{p_n})$ of \cref{csf:lem:colors} in place of an abstract family, as the manuscript itself says (\cref{csf:sec:intersection}, and after \cref{csf:thm:universality}); the unrestricted theorem stays credited to \cite{cmms}. Rast--Sahota, Miller's lecture notes, Marcone--Rosendal and Laver were not rechecked.

\emph{Neighbouring reports.} Detailed notes follow the results concerned.
\begin{itemize}[leftmargin=1.5em]
\item \repo{foundations-and-computation/definable-surreals-and-omnific-integers} and Part~VII of \repo{foundations-and-computation/birthday-cutoffs-and-hereditary-sets}: \cref{csf:lem:independence} is a case of the former's Lemma~8.1 (\lbl{dsn:lem:cosets}), and an analogue of the latter's \lbl{hset:kw:thm:double}; see the note after \cref{csf:lem:independence}. Part~VII otherwise proves different theorems (choiceless coding and universality of $\No$ over ZF), and shares no other result with this report.
\item \repo{surreal/surreal-self-embeddings}, cited as background \cite{sse}: \cref{csf:lem:valuation-functor} is the subfield form of its \lbl{sse:prop:valueaction} (note after the lemma). Its questions concern self-embeddings of $\No$, not countable subfields; none is answered here.
\item \repo{surreal/independent-surreal-copies}: its \lbl{isc:thm:boolean} realizes a Boolean lattice of full Hahn fields, the analogue of \cref{csf:cor:lattice} (note after the corollary).
\item \repo{foundations-and-computation/polish-models-of-omnific-arithmetic}, Part~II: its \lbl{pma:bor:lem:KB} contains \cref{csf:lem:kb} and the analytic-completeness assertion of \cref{csf:thm:fixed-hard} in complement form (note after \cref{csf:thm:fixed-hard}); its \lbl{pma:bor:thm:dichotomy} bears on \cref{csf:sec:questions}. Part~XV (batch~92): its \lbl{pma:rk:thm:scaling} applies the integer comparisons of \cref{csf:lem:valuation-functor} to left-finite Levi-Civita fields (note after the lemma).
\item \repo{foundations-and-computation/surreal-fields-across-universes}, Part~VIII: \lbl{univ:gs:thm:boolean} also indexes a family of real closed fields by subsets of a set, the class fields $\No^{M_A}$ of Cohen extensions, but there embeddability is exactly inclusion, $A\subseteq A'$; here it is embeddability of colored rank orders (\cref{csf:thm:colored-coding}). Different objects, no shared proof.
\item \repo{foundations-and-computation/omnific-notations} (Part~IV: $\boldsymbol\Pi^1_1$-complete support validity) uses the same ill-founded-tree toolkit on different objects; \repo{foundations-and-computation/computable-surreals} constructs countable real closed subfields of $\No$ with effective presentations (its main theorem, label \texttt{thm:main} there), background for Question~\ref{csf:q:effective} but no shared theorem.
\item Batch-90 siblings, written concurrently: manuscript~06 (source~04 of \repo{surreal/discrete-initial-subgroups-and-omnific-normalization}, labels \texttt{isg:ptg:}) also cites \cite{glazer}; manuscripts 01 and~03 became Parts~XI--XII of \repo{polish-models-of-omnific-arithmetic}, 05 Part~VIII of \repo{birthday-cutoffs-and-hereditary-sets}, and 02 Part~II of the research-report collection's \repo{ordinals-and-order-types/measurable-box-games}. None shares a theorem with this report.
\end{itemize}

\emph{The write.} Being a single source, the report needed no merge choices. The write prefixed all 56 delivered labels with \texttt{csf:} and updated every reference (63 \verb|\cref| keys and 15 \verb|\eqref|); added the labels \texttt{csf:sec:provenance}, \texttt{csf:sec:notation} and twelve question labels \texttt{csf:q:*} (the delivered questions had none); added these dated \textbf{[write]} notes, a title-page line and the \texttt{pageanchor} switch around the title page; and gave the labels of the 8 lemmas, 4 corollaries, 2 propositions and 12 questions cleveref type hints (such as \verb|\label[lemma]|), because these environments share the theorem counter and the delivered text's references named every one of them ``theorem''. No statement, proof, number or non-claim of the manuscript was changed, and no section, theorem, equation or question number moved.
\end{writenote}

\subsection{Notation}
\label{csf:sec:notation}

\begin{writenote}[Added 4 October 2026, batch~90]
The manuscript's symbols, each with its meaning here and the reading to avoid. No symbol was renamed.
{\footnotesize
\begin{longtable}{P{0.2\textwidth}P{0.4\textwidth}P{0.3\textwidth}}
\toprule
Symbol & Meaning here & Not to be read as\\
\midrule
\endhead
$\omega^x$, $\e{x}$ & Conway's omega-map (monomials) & Gonshor's $\exp$\\
$k$ & the field of real algebraic numbers, fixed & a variable field\\
$p_n$, $B$, $I$ & the $n$th prime; $\{1\}\cup\{\sqrt{p_n}\}$; $\Q\times B$ & \\
$b$ & an element of $B$ (a real coefficient) & the birthday $b(a)$ of \cref{csf:sec:foundations}\\
$x_{q,b}$, $t_q$ & $\omega^{b\omega^q}$; $t_q=x_{q,1}$ & \\
$K_A$, $\Kstar$, $\Kone_S$, $\Kone$ & relative real closures (\cref{csf:eq:ambient}); $\Kone_S=K_{S\times\{1\}}$ & $K_A$ is not a Hahn field\\
$\Cstar$, $\Cone$ & the families $\{K_A\}$, $\{\Kone_S\}$, subsets of $2^{\Kstar}$ & \\
$\rc_{\No}$ & relative real closure inside $\No$ (a set) & a class construction\\
$D(a)$ & finite generator support (\cref{csf:thm:support}) & normal-form support $\supp$\\
$D$, $P_-$ & the order $\{-n\}$ of type $\omega^*$; $P_-=\Kone_D$ & $D(a)$, $D_A$\\
$D_A$, $A_q$ & complete-column ranks $\{q:A_q=B\}$; the column $\{b:(q,b)\in A\}$ & \\
$J_A$ & $\{q:H_A(q)\ne0\}$ (\cref{csf:thm:valuation}) $=\{q:A_q\ne\varnothing\}$ (\cref{csf:sec:generic-star}); the two definitions agree & \\
$H_A(q)$, $\Gamma_A$ & component $\operatorname{span}_\Q\{b:(q,b)\in A\}$; value group (\cref{csf:eq:value-group}) & \\
$H_n$, $H$ & $\Q(\sqrt{p_n})$ as an ordered group; $\operatorname{span}_\Q B$ & ordered additive groups; $H_n$ is also a field (used in \cref{csf:lem:colors}), $H$ is not\\
$v$ & natural valuation, $v(x)=-(\text{leading exponent})$ & a Conway exponent\\
$E^\star_P$, $E^{(1)}_P$ & fixed-source embedding loci, sets of parameters & a predicate or relation\\
$\Sig$, $\PiC$ & boldface analytic, coanalytic & a $\sigma$-algebra\\
analytic-complete & a set to which every analytic set reduces continuously & complete analytic \emph{quasiorder} (\cref{csf:sec:descriptive})\\
$\KB$, $\IF$, $\WF$, $\mathsf{Tr}$ & Kleene--Brouwer order; ill-founded and well-founded trees; trees & \\
$S_L$, $e_L$ & rational placement of an order $L$ on $\N$ & $S_T$ (tree code)\\
$S_T$, $W$, $h$ & tree code (\cref{csf:eq:tree-code}) and its prefix; $h$ also names other order embeddings & \\
$Z(L)$, $G$, $\psi$, $R_T$, $\mathcal N$ & integer-sum transform; $G=Z(\KB)$; $\mathcal N$ non-scattered suborders of $\Q$ & \\
$M_L$, $A_L$ & colored-order code (\cref{csf:eq:color-code}) & \\
$\mathcal G$, $\mathcal W_P$, $\tau_r$, $\mu$ & sets $S$ dense in $\Q$; witness relation; translation; coin-flip measure & a measure on $\No$\\
$\beta$, $b(a)$ & a birthday cutoff and the birthday of $a$ & an optimal bound\\
\bottomrule
\end{longtable}
\par}
Across the collection: $K_A$ is not \repo{independent-surreal-copies}' $E_J$ (full Hahn fields), Part~VII's $T_X$ or $U_\kappa$ of \repo{birthday-cutoffs-and-hereditary-sets}, or \repo{surreal-fields-across-universes}' $\No^{M_A}$; the tree-coding embedding $h$ plays the role of $\mu_{\KB}$ in \repo{polish-models-of-omnific-arithmetic}; and $E^\star_P$ is unrelated to the enumeration predicate $E$ of Parts~VI and~VIII of \repo{birthday-cutoffs-and-hereditary-sets}.
\end{writenote}

\section{Foundations, conventions, and the proof boundary}
\label{csf:sec:foundations}

\subsection{Set-sized relative real closures}

We use the standard surreal ordered field and its normal-form theorem
\cite{gonshor,vde}. For a set-sized ordered subfield $F\subset\No$, define
\[
 \rc_{\No}(F)=\{a\in\No:a\text{ is algebraic over }F\}.
\]
Although the notation quantifies over $\No$, its output is a set. Indeed,
$F[X]$ is a set, and a nonzero polynomial has at most its degree many roots
in any field. For each polynomial and each possible root index, its ordered
list of roots, when present in $\No$, is uniquely determined. Replacement
collects these roots into a set. Their union is exactly the displayed class.

This set is an algebraic extension field of $F$. It is real closed: positive
square roots and roots of odd-degree polynomials exist in $\No$, and any
such root of a polynomial with coefficients algebraic over $F$ is again
algebraic over $F$. When $F$ is countable, there are only countably many
polynomials and finitely many roots per nonzero polynomial, so the relative
real closure is countable.

Equivalently, one can first construct the abstract real closure of $F$ and
use its unique order-preserving embedding into $\No$ over $F$. Once
$\Kstar$ has been constructed, every subsequent definition can be written
entirely inside the set $\Kstar$:
\begin{equation}
 K_A=\{a\in\Kstar:a\text{ is algebraic over }k(x_i:i\in A)\}.
 \label{csf:eq:set-relative}
\end{equation}
The two definitions agree. In particular, $\Cstar$ is a set of subsets of a
countable set, not a collection requiring quantification over arbitrary
proper classes.

\subsection{One ordinal cutoff suffices}

Let $b(a)$ denote the birthday of a surreal $a$. Since $\Kstar$ is a set,
\[
 \beta=\sup\{b(a)+1:a\in\Kstar\}
\]
is an ordinal, and $\Kstar\subseteq\No_{<\beta}$. Therefore the same birthday
cutoff contains every field in both families. No numerical bound for
$\beta$ is needed or asserted. The cutoff itself is not being asserted to
be a field, and its full structure is not used.

This is also a useful type-theoretic separation. A formalization can work
with a set-sized ambient real closed field, its countable independent family,
and relative algebraic closures. The embedding into a suitable universe of
surreals is a separate interface. One need not form a type of all classes
of surreal numbers.

\subsection{Embeddings and elementarity}

By an embedding we mean a unital field homomorphism. Such a map between real
closed fields is injective and order preserving: positive elements are
exactly the nonzero squares. It reflects the order as well. It fixes $k$
pointwise, because a real algebraic number is specified uniquely by its
minimal polynomial over $\Q$ and an isolating interval with rational
endpoints.

We use two standard consequences of real-closed-field theory. An ordered
field has a real closure unique up to a unique ordered isomorphism over that
field; and every embedding between real closed fields is elementary in the
language of ordered rings, by quantifier elimination for $\RCF$
\cite{marker}. Thus ``elementary'' in \cref{csf:thm:main} introduces no additional
coding predicate or saturation assumption.

\subsection{External inputs versus new proofs}

The external inputs are surreal real closedness and normal forms, elementary
field-theoretic algebraic dependence, real-closure uniqueness, quantifier
elimination for $\RCF$, completeness of ill-founded trees, Borel completeness
of countable-order isomorphism, and completeness of colored-order
embeddability. All other steps needed for the stated package are proved
below. In particular, neither an unproved universality assertion about
$\No$ nor a global extension of our embeddings to $\No$ is needed.

\section{Monomial coordinates and their exact valuation invariants}
\label{csf:sec:coordinates}

\subsection{Independent algebraic real coefficients}

For a finite set of distinct primes $p_1,\ldots,p_m$, the real multiquadratic
field $\Q(\sqrt{p_1},\ldots,\sqrt{p_m})$ has basis
\[
 \left\{\sqrt{\prod_{j\in J}p_j}:J\subseteq\{1,\ldots,m\}\right\}.
\]
Here is an elementary induction justifying the basis and all independent
sign changes. Assume the statement for the first $m$ primes, and suppose
$\alpha=\sqrt{p_{m+1}}$ belonged to the field they generate. Each of its
sign-changing automorphisms sends $\alpha$ to $\alpha$ or $-\alpha$,
since these are the roots of $X^2-p_{m+1}$. In the asserted basis for the
first $m$ primes, simultaneous eigenvectors of all sign changes are
rational multiples of individual basis monomials: distinct monomials have
distinct sign characters. Thus
$\alpha=c\sqrt{\prod_{j\in J}p_j}$ for some $c\in\Q$.
Squaring contradicts the odd valuation at the new prime $p_{m+1}$.
The next adjunction therefore has degree two. Old sign changes extend
while fixing the new square root, and the new square root has its own
sign change. This proves the induction, beginning with $\Q$.

\begin{lemma}
\label[lemma]{csf:lem:radicals}
The set $B$ in \eqref{csf:eq:ambient-generators} is linearly independent over
$\Q$. More generally, distinct squarefree radical monomials in finitely
many primes are $\Q$-linearly independent.
\end{lemma}
\begin{proof}
Each alleged dependence lies in one finite multiquadratic field. Apply its
independent sign-changing automorphisms, multiply by the corresponding
characters, and average. Each coefficient is isolated and hence zero.
\end{proof}

For a reader wishing to avoid the degree formulation, the same elementary
multiquadratic lemma is a standard finite algebraic input; the verification
script includes an exact character-orthogonality check, but that finite
check is not the proof for arbitrary $m$.

\begin{writenote}[Added 4 October 2026, batch~90]
For arbitrary $m$, and indeed for arbitrarily many primes, this lemma is formally proved in ProveIt's Lean development: it is the instance $M=\Q$, $a_i=p_i$, with square roots in $\R$, of \texttt{Surreal.TailSigns.linearIndependent\_mono} (\lbl{tail:lem:signs} of \repo{surreal/tail-spans-and-differential-transcendence}, proved by the same sign-change induction), whose square-class hypothesis for distinct primes is \texttt{Surreal.AlgebraicClosurePart.independentSquareClasses\_prime}. The specialization is not itself a declaration; see \cref{csf:sec:provenance}.
\end{writenote}

\subsection{The transcendence basis}

A finite surreal normal form
\[
 g=\sum_{q\in F} r_q\omega^q,\qquad r_q\in\R,
\]
is positive exactly when the coefficient at the largest $q$ with
$r_q\ne0$ is positive. For distinct exponents $g<h$, the ratio
$\omega^h/\omega^g=\omega^{h-g}$ exceeds every positive real number.
Also $\omega^{g+h}=\omega^g\omega^h$.

\begin{lemma}[Independent double monomials]
\label[lemma]{csf:lem:independence}
The family $(x_{q,b})_{(q,b)\in I}$ is algebraically independent over $k$.
\end{lemma}
\begin{proof}
An ordinary monomial with finite nonnegative integer exponent vector
$\nu=(\nu_{q,b})$ has value
\[
 \prod_{q,b}x_{q,b}^{\nu_{q,b}}
 =\omega^{g_\nu},\qquad
 g_\nu=\sum_q\left(\sum_b\nu_{q,b}b\right)\omega^q.
\]
If $\nu\ne\mu$, then $g_\nu\ne g_\mu$ by \cref{csf:lem:radicals} and uniqueness
of the finite normal form. A nonzero polynomial in finitely many $x_{q,b}$
therefore becomes a finite sum of distinct surreal monomials with nonzero
coefficients in $k$. Its largest-exponent summand dominates all other
summands, so the sum is nonzero.
\end{proof}

Consequently
\begin{equation}
 \trdeg_k K_A=|A|,
 \label{csf:eq:trdeg}
\end{equation}
where the right side is a finite cardinal or $\aleph_0$. The use of real
algebraic $b$ does not make $x_{q,b}$ algebraic over $x_{q,1}$: an irrational
multiple of the exponent is precisely what prevents an integral polynomial
relation.

\begin{writenote}[Added 4 October 2026, batch~90]
\Cref{csf:lem:independence} is a case of Lemma~8.1 (\lbl{dsn:lem:cosets}, ``Disjoint support cosets'', that report's source~04) of \repo{foundations-and-computation/definable-surreals-and-omnific-integers}: take the subfield $k$, whose nonzero elements have normal-form support $\{0\}$, the subgroup $G=\{0\}$, and $g_{q,b}=b\omega^q$; a nonzero finitely supported integer combination of the $g_{q,b}$ is nonzero by \cref{csf:lem:radicals} and uniqueness of the finite normal form, so the monomials $\omega^{g_{q,b}}=x_{q,b}$ are algebraically independent over $k$. Part~VII of \repo{foundations-and-computation/birthday-cutoffs-and-hereditary-sets} proves the coefficient-one analogue over $\R$, the double monomials $\omega^{\omega^{j(x)}}$ (Theorem~43.1, \lbl{hset:kw:thm:double}), by the same leading-term argument, and its Remark \lbl{hset:kw:rem:other} records that it too is a case of \lbl{dsn:lem:cosets}. Both are credited; the proof above is printed as the manuscript's own and claims no novelty beyond the radical coefficients $b\in B$, which make several independent generators share one rank.
\end{writenote}

\subsection{Natural valuation}

For an ordered field $F$, its natural valuation ring and maximal ideal are
\[
 \mathcal O_F=\{a:|a|\le n\text{ for some }n\in\N_{>0}\},\qquad
 \mathfrak m_F=\{a:|a|<1/n\text{ for every }n\in\N_{>0}\}.
\]
The value group is the multiplicative group modulo bounded positive units,
written additively with the usual valuation orientation. For a nonzero
surreal with leading normal-form exponent $g$, we identify its value with
$-g$. Thus $v(x_{q,b})=-b\omega^q$.

For $A\subseteq I$, set
\begin{equation}
 H_A(q)=\operatorname{span}_{\Q}\{b:(q,b)\in A\}\subseteq\R,
 \qquad
 \Gamma_A=\bigoplus_{q\in\Q}H_A(q)\omega^q.
 \label{csf:eq:value-group}
\end{equation}
The direct sum consists of finite-support vectors. It is ordered by the
largest nonzero coordinate and its real coefficient. Empty components are
omitted from its rank chain.

\begin{theorem}[Exact value group and residue field]
\label{csf:thm:valuation}
For every $A\subseteq I$, the natural value group of $K_A$ is $\Gamma_A$,
and its residue field is exactly $k$. If
$J_A=\{q:H_A(q)\ne0\}$, then the Archimedean rank chain of $\Gamma_A$ is
$(J_A,<)$, with component $H_A(q)$ at rank $q$.
\end{theorem}
\begin{proof}
Write $F_A=k(x_i:i\in A)$. The leading-exponent computation shows that
$v(F_A^\times)$ is the integral span of the $b\omega^q$ for $(q,b)\in A$.
Every rational multiple of such a value occurs in $K_A$, by taking positive
integer roots of monomials. Hence $\Gamma_A\subseteq v(K_A^\times)$.

Conversely, let $a\in K_A^\times$ satisfy a nonzero polynomial relation
$\sum_{i=0}^d c_i a^i=0$ over $F_A$. Among its nonzero summands, the minimum
of $v(c_i)+iv(a)$ is attained at least twice: a unique least-value term
cannot cancel with terms of greater value. For two distinct such indices,
\[
 (i-j)v(a)=v(c_j)-v(c_i).
\]
The right side lies in the integral span just computed. The ordered value
group is torsion free, so $v(a)$ lies in its rational span $\Gamma_A$.
This proves the value-group assertion.

The residue of a unit in $F_A$ is the ratio of two leading coefficients,
and therefore lies in $k$. For a unit $a\in K_A$, choose a polynomial
relation over $F_A$ and divide all coefficients by a coefficient of least
value. All normalized coefficients belong to the valuation ring of $F_A$,
and at least one has nonzero residue. Reduction yields a nonzero polynomial
over $k$ satisfied by the residue of $a$. The residue field is ordered and
$k$ is real closed, so this algebraic residue belongs to $k$. Elements of
positive value have residue zero, and all constants in $k$ occur. Thus the
residue field is $k$.

For the rank assertion, two positive elements of $\Gamma_A$ are
Archimedean equivalent exactly when their largest nonzero coordinates
coincide. If the largest coordinates differ, no integer multiple of the
smaller-rank element reaches the larger-rank one. At a fixed coordinate,
the leading coefficients are positive real numbers, so sufficiently large
integer multiples compare them in both directions. Finally, quotienting
vectors supported at coordinates $\le q$ by those supported at coordinates
$<q$ gives exactly $H_A(q)$.
\end{proof}

\subsection{Why arbitrary field embeddings preserve these invariants}

\begin{lemma}[Intrinsic value-group functor]
\label[lemma]{csf:lem:valuation-functor}
An embedding $f:F\emb G$ of ordered fields induces an ordered-group
embedding $\bar f:v(F^\times)\emb v(G^\times)$. It also induces an order
embedding of Archimedean rank chains and, at each source rank, an ordered
embedding of the corresponding Archimedean components.
\end{lemma}
\begin{proof}
All comparisons $|a|\le n|b|$ are preserved and reflected by $f$. Hence
bounded-ratio equivalence and the comparison of natural values are
preserved and reflected. The resulting map on values is well defined,
injective, and additive because $f$ respects multiplication.

For positive value-group elements, Archimedean equivalence is expressed by
two inequalities of the form $g\le nh$. These too are preserved and
reflected. At the rank of a positive $g$, form the convex subgroups
\[
 C_g=\{h:|h|\le ng\text{ for some }n\ge1\},\qquad
 D_g=\{h:n|h|<g\text{ for every }n\ge1\}.
\]
The map sends $C_g$ into $C_{\bar f(g)}$ and sends $D_g$ into
$D_{\bar f(g)}$, with the latter membership reflected. It therefore induces
an injective ordered homomorphism between the quotients.
\end{proof}

This lemma is essential. Restricting attention to generator-preserving
embeddings would prove only a sufficient condition. The intrinsic
valuation recovers a necessary condition for \emph{every} field embedding.

\begin{writenote}[Added 4 October 2026, batch~90]
The first assertion is the subfield form of \lbl{sse:prop:valueaction} of \repo{surreal/surreal-self-embeddings}, which proves by the same integer-comparison argument that every field self-embedding $j$ of $\No$ induces an ordered value-group embedding $\tau_j$; the component maps on convex quotients are the step used in the proof of Proposition~7.2 of \cite{cmms} (see \cref{csf:sec:provenance}).
\end{writenote}

\begin{writenote}[Added 4 October 2026, batch~92]
Part~XV of \repo{foundations-and-computation/polish-models-of-omnific-arithmetic} (batch-92 manuscript~05; AI-assisted, unrefereed, not formalized) uses the same integer comparisons for field embeddings $\mathcal L^\Gamma\to\mathcal L^\Delta$ between left-finite Levi-Civita fields with real algebraic coefficients and exponent groups $0\ne\Gamma,\Delta\le\R$, where the induced value-group map is multiplication by one real $\lambda>0$ (\lbl{pma:rk:thm:scaling}), and it classifies these embeddings by triples $(\lambda,c,u)$ (\lbl{pma:rk:thm:parameters}). Those fields are complete and uncountable, not countable subfields of $\No$, so \cref{csf:q:embeddings} is not answered. Its Cantor family of closed immediate self-copies of one Polish field (\lbl{pma:rk:thm:family}) is a different object from the Cantor families of countable subfields here.
\end{writenote}

\section{Finite supports and a compact Boolean family}
\label{csf:sec:support}

\subsection{A general algebraic lemma}

The next theorem is not specific to surreal numbers. Let $F$ be a field,
let $(t_i)_{i\in I}$ be algebraically independent over $F$, and suppose
$M$ is algebraic over $F(t_i:i\in I)$. For $A\subseteq I$, define
\[
 M_A=\{a\in M:a\text{ is algebraic over }F(t_i:i\in A)\}.
\]

\begin{theorem}[Unique finite support]
\label{csf:thm:support}
For every $a\in M$ there exists a unique finite $D(a)\subseteq I$ such that
\begin{equation}
 a\in M_A\quad\Longleftrightarrow\quad D(a)\subseteq A
 \qquad(A\subseteq I).
 \label{csf:eq:support}
\end{equation}
Moreover $D(t_i)=\{i\}$, and $D(a)=\varnothing$ exactly when $a$ is algebraic
over $F$.
\end{theorem}
\begin{proof}
A polynomial witnessing that $a$ is algebraic over $F(t_i:i\in I)$ uses
only finitely many generators. Choose a finite set $D$ minimal under
inclusion such that $a$ is algebraic over $F(t_i:i\in D)$.

We first recall the elementary exchange step. If $a$ is algebraic over
$L(t)$ but transcendental over $L$, clearing denominators gives a nonzero
$P(X,Y)\in L[X,Y]$ with $P(a,t)=0$. Viewed as a polynomial in $Y$, it
remains nonzero after $X=a$, because $a$ is transcendental over $L$.
Thus $t$ is algebraic over $L(a)$.

Now suppose $a\in M_A$. Some finite $C\subseteq A$ already suffices for
this algebraicity. If $i\in D\setminus C$, minimality of $D$ and exchange
show that $t_i$ is algebraic over
$F(a,t_j:j\in D\setminus\{i\})$. Since $a$ is algebraic over $F(t_j:j\in C)$,
transitivity of algebraicity makes $t_i$ algebraic over
\[
 F(t_j:j\in C\cup(D\setminus\{i\})).
\]
This contradicts algebraic independence, because $i$ is absent from the
index set. Consequently $D\subseteq C\subseteq A$.

The converse follows by monotonicity of algebraicity. This proves
\eqref{csf:eq:support}; it also proves uniqueness. The final two assertions
follow immediately from independence and the definition.
\end{proof}

\begin{remark}[Generator support is not normal-form support]
An algebraic element can have infinitely many terms in its surreal normal
form and still have a one-element generator support. For example,
$\sqrt{1+x_i^{-1}}$ is algebraic over $k(x_i)$ and its support in the present
sense is $\{i\}$, even though its usual infinitesimal expansion has
infinitely many terms. No claim that algebraic normal forms are finite is
being made.
\end{remark}

\subsection{Topological consequences}

Apply \cref{csf:thm:support} to $F=k$, $t_i=x_i$, and $M=\Kstar$. Fix an
enumeration of the countably infinite set $I$ and an enumeration of
$\Kstar$. The power sets $2^I$ and $2^{\Kstar}$ carry their ordinary product
Cantor topologies.

\begin{theorem}[Closed coordinate family]
\label{csf:thm:cantor}
The map $\Phi(A)=K_A$ is continuous and injective, hence a homeomorphism
onto a closed subspace $\Cstar\subseteq2^{\Kstar}$ homeomorphic to Cantor
space. Every coordinate test $a\in K_A$ is clopen as a condition on $A$.
Every $K_A$ is an elementary real closed subfield of $\Kstar$.
\end{theorem}
\begin{proof}
For a fixed $a\in\Kstar$, \cref{csf:thm:support} gives
\[
 \{A:a\in K_A\}=\bigcap_{i\in D(a)}\{A:i\in A\},
\]
a finite intersection of clopen sets. Thus every membership coordinate of
$\Phi$ is continuous. Since $x_i\in K_A$ exactly when $i\in A$, the map is
injective. A continuous injection from compact $2^I$ into Hausdorff
$2^{\Kstar}$ is a homeomorphism onto its compact, and hence closed, image.
Real closedness was established in \cref{csf:sec:foundations}; elementarity
follows from quantifier elimination.
\end{proof}

\begin{corollary}[Coordinate lattice]
\label[corollary]{csf:cor:lattice}
For arbitrary families $(A_j)_{j\in J}$,
\[
 \bigcap_{j\in J}K_{A_j}=K_{\bigcap_j A_j},\qquad
 \rc_{\Kstar}\left(\text{field generated by }\bigcup_jK_{A_j}\right)
 =K_{\bigcup_jA_j}.
\]
For a nonempty directed family under inclusion,
\[
 \bigcup_jK_{A_j}=K_{\bigcup_jA_j}.
\]
Consequently $\Cstar$, with these relative meet and join operations, is a
complete Boolean lattice isomorphic to $\mathcal P(I)$.
\end{corollary}
\begin{proof}
The intersection assertion is \eqref{csf:eq:support}. For the join, every
member on the left is algebraic over the field generated by coordinates
in the union, while that field contains every required generator. Relative
real closure gives equality. In the directed case, the finite support of
an element in the right-hand side is already contained in one member of
the family. The Boolean complement of $K_A$ within this coordinate lattice
is $K_{I\setminus A}$; the lattice bottom is $k$.
\end{proof}

These joins are not literal unions of two fields. For instance,
$x_i+x_j$ can lie in $K_{\{i,j\}}$ without belonging to either
$K_{\{i\}}$ or $K_{\{j\}}$.

\begin{writenote}[Added 4 October 2026, batch~90]
\repo{surreal/independent-surreal-copies} proves a Boolean lattice of the same shape for \emph{full Hahn fields}: for an independent family of exponent groups $G_i$ over a core $H$, the fields $E_J$ on $G_J$ satisfy $E_J\cap E_T=E_{J\cap T}$, with full Hahn join $E_{J\cup T}$ (\lbl{isc:thm:boolean}). There the fields are set-sized or class-sized Hahn fields and the proof uses normal-form supports; here they are countable relative real closures and the proof uses algebraic exchange (\cref{csf:thm:support}). No result is shared.
\end{writenote}

\subsection{Uniform codes on the domain \texorpdfstring{$\N$}{N}}

The family can be transferred to the usual space of countable structures
without adding complexity. Fix an enumeration $(z_n)$ of $\Kstar$ without
repetitions, and enumerate each $K_A$ by the increasing indices of its
members in this list. The $m$th enumerated element depends on only finitely
many tests $z_n\in K_A$, and hence continuously on $A$. The inherited
addition, multiplication, order, and constants can then be coded on $\N$.
Every individual operation-table entry depends locally on finitely many
membership tests. This yields a continuous presentation map into the
space of structures in the countable ordered-field language.

Continuity here is relative to a fixed enumeration and the product
structure space. It does not assert an effective algorithm for computing
all surreal algebraic numbers or their supports from an arbitrary external
notation.

\section{The rank-one family remembers exactly a countable order}
\label{csf:sec:rank-one}

The rank-one field $\Kone_S$ is generated, up to real closure, by
$t_q=x_{q,1}=\omega^{\omega^q}$ for $q\in S$. Its natural value group is
\begin{equation}
 \Gamma^{(1)}_S=\bigoplus_{q\in S}\Q\omega^q.
 \label{csf:eq:rank-one-values}
\end{equation}

\begin{theorem}[Exact embedding and isomorphism classification]
\label{csf:thm:rank-one-classification}
For arbitrary $S,T\subseteq\Q$,
\begin{align}
 \Kone_S\emb\Kone_T
 &\quad\Longleftrightarrow\quad (S,<)\emb(T,<),
 \label{csf:eq:embedding-classification}\\
 \Kone_S\cong\Kone_T
 &\quad\Longleftrightarrow\quad (S,<)\cong(T,<).
 \label{csf:eq:iso-classification}
\end{align}
The assertions include finite and empty index sets.
\end{theorem}
\begin{proof}
An embedding of fields induces an embedding of natural value groups by
\cref{csf:lem:valuation-functor}. The Archimedean rank chains of
\eqref{csf:eq:rank-one-values} are $S$ and $T$, respectively, so this induces
an order embedding $S\emb T$. An isomorphism induces an order isomorphism
of ranks.

Conversely, given an order embedding $h:S\emb T$, send $t_q$ to $t_{h(q)}$.
Algebraic independence makes this a field embedding on the rational-function
fields. It preserves order: for any finite polynomial, the comparison of
two monomial exponents is determined by the largest coordinate at which
their integer exponent vectors differ. That coordinate and its coefficient
are preserved by $h$. The sign of the leading coefficient is unchanged.
The same follows for quotients.

The ordered embedding extends uniquely to real closures. Its image is
$\Kone_{h(S)}\subseteq\Kone_T$. If $h$ is onto, the extension is onto.
For $S=\varnothing$, the source is $k$ and its value rank is empty, as
required.
\end{proof}

\begin{example}[Equal coarse invariants, different fields]
Let $S$ have order type $\omega$ and $T$ order type $\omega^*$.
Both fields have residue field $k$, countably infinite transcendence degree,
and value groups abstractly isomorphic to the $\Q$-vector space
$\Q^{(\N)}$. As bare ordered sets, both fields have order type $\Q$.
Nevertheless they are not isomorphic, and neither embeds into the other,
because neither of these two rank orders embeds into the other.
\end{example}

\begin{remark}[No classification of all individual maps]
\cref{csf:thm:rank-one-classification} classifies the \emph{existence} of an
embedding and the resulting isomorphism types. It does not say that every
embedding sends each displayed monomial to a displayed monomial. An
embedding can contain additional unit factors or other algebraic choices.
Only its induced rank map is used in the necessary direction.
\end{remark}

\section{Descriptive complexity in the rank-one family}
\label{csf:sec:descriptive}

\subsection{Three different completeness notions}

An analytic set is a projection of a Borel set in a product of standard
Borel spaces. An analytic subset of a Polish space is
\emph{analytic-complete} if every analytic subset of Baire space continuously
reduces to it. The completeness statements for individual sets below use
this convention. We exhibit explicit continuous reductions from the space
of ill-founded trees.

For quasiorders, the stronger phrase \emph{complete analytic quasiorder}
means that for every analytic quasiorder $R$ on a standard Borel space $X$
there is a Borel map $f$ with
\[
 x\,R\,y\quad\Longleftrightarrow\quad f(x)\,Q\,f(y).
\]
Both coordinates use the same map. Analytic-completeness of a relation
viewed merely as a set of pairs does not imply this universal property.
Finally, \emph{Borel complete for countable structures}, or
$S_\infty$-complete, refers to reducing arbitrary isomorphism relations on
countable structures to the target isomorphism relation.

\subsection{A Borel rational placement of countable orders}

Let $L$ be a strict linear order on $\N$. Choose once and for all an
enumeration of $\Q$. Recursively select $e_L(n)$ to be the first unused
rational satisfying
\[
 e_L(i)<e_L(n)\quad\Longleftrightarrow\quad i<_L n
 \qquad(i<n).
\]
At each stage, only finitely many inequalities are imposed. Density and
lack of endpoints of $\Q$ ensure a solution. The first $n+1$ choices depend
only on the restriction of $L$ to $\{0,\ldots,n\}$, so each coordinate
$e_L(n)$ is Borel, indeed locally constant. Therefore
\[
 S_L=\{e_L(n):n\in\N\}
\]
is a Borel function of $L$, because $q\in S_L$ is a countable union of the
Borel conditions $e_L(n)=q$.

\begin{corollary}[Borel-complete isomorphism]
\label[corollary]{csf:cor:isomorphism}
Isomorphism on $\Cone$, and hence on $\Cstar$, is Borel complete for
countable structures.
\end{corollary}
\begin{proof}
The map $L\mapsto\Kone_{S_L}$ is a Borel reduction of countable-linear-order
isomorphism by \cref{csf:thm:rank-one-classification}. The latter isomorphism
relation is Borel complete \cite{friedman,rast}. The upper classification
bound follows because our fields are countable structures in a countable
language, with the uniform codes described above.
\end{proof}

\subsection{A fixed source with an analytic-complete locus}

Let $\mathsf{Tr}$ be the closed subspace of $2^{\N^{<\N}}$ consisting of
prefix-closed trees. Its members may be empty. A branch is a function
$f\in\N^\N$ whose finite initial segments all belong to the tree. The set
$\IF$ of trees with a branch is analytic-complete
\cite[Proposition 1.4.28]{miller}. Its complement $\WF$ is
coanalytic-complete.

Fix the Kleene--Brouwer order on $\N^{<\N}$: a proper extension precedes
its initial segment, and otherwise the first unequal entries are compared
in their usual order. Write it as $<_{\KB}$.

\begin{lemma}
\label[lemma]{csf:lem:kb}
A tree $T$ has a branch exactly when its Kleene--Brouwer order is not a
well-order.
\end{lemma}
\begin{proof}
The successive initial segments of a branch form a strictly decreasing
Kleene--Brouwer sequence. Conversely, if $T$ has no branch, the relation
of proper extension on its nodes is well founded. By well-founded
induction, the Kleene--Brouwer order on the subtree above a node $s$ is
a well-order: it is the ordinal sum, in increasing order of $n$, of the
orders on the subtrees above $s^\frown n$, followed by the node $s$ itself.
Every summand is a well-order by the induction hypothesis, and a sum of
well-orders over a subset of $\N$ is a well-order. Apply this at the root;
the empty-tree case is immediate.
\end{proof}

Choose one order embedding
\[
 h:\omega+(\N^{<\N},<_{\KB})\emb\Q.
\]
It exists because this is a countable linear order. Denote the image of
the first summand by $W$, and for $T\in\mathsf{Tr}$ put
\begin{equation}
 S_T=W\cup\{h(s):s\in T\}.
 \label{csf:eq:tree-code}
\end{equation}
This map is continuous: membership at a fixed rational coordinate is
constant, or is one membership bit of $T$.

\begin{theorem}[One fixed field already has complete analytic existence]
\label{csf:thm:fixed-hard}
Let $P_-=\Kone_D$, where $D$ has order type $\omega^*$. Then
\[
 E^{(1)}_{P_-}=\{S:(S,<)\text{ is not a well-order}\}
\]
is analytic-complete. Its complement is coanalytic-complete. Moreover,
\[
 T\in\IF\quad\Longleftrightarrow\quad P_-\emb\Kone_{S_T},
\]
where every field $\Kone_{S_T}$ has transcendence degree $\aleph_0$ over
$k$, residue field $k$, and abstract value group $\Q^{(\N)}$.
\end{theorem}
\begin{proof}
By \cref{csf:thm:rank-one-classification}, $P_-\emb\Kone_S$ exactly when
$\omega^*$ embeds into $S$, that is, exactly when $S$ admits an infinite
strictly decreasing sequence. For countable linear orders in ZFC, this is
precisely failure of well-foundedness. The existence of such a sequence is
analytic.

In \eqref{csf:eq:tree-code}, the added $\omega$ summand is a well-order and
cannot create an infinite decreasing sequence. Therefore the target order
is ill founded exactly when the Kleene--Brouwer order on $T$ is ill
founded, which is equivalent to $T\in\IF$ by \cref{csf:lem:kb}. The reduction is
continuous. The fixed infinite prefix $W$ ensures $|S_T|=\aleph_0$, even
for empty or finite $T$. The algebraic invariants follow from
\cref{csf:thm:valuation} and \eqref{csf:eq:trdeg}.
\end{proof}

\begin{writenote}[Added 4 October 2026, batch~90]
Part~II of \repo{foundations-and-computation/polish-models-of-omnific-arithmetic} proves the same Kleene--Brouwer lemma inside its Lemma~\lbl{pma:bor:lem:KB}, with a fixed order embedding $\mu_{\KB}:\N^{<\N}\to\Q\cap(0,1)$ in the role of $h$, and concludes that $\mathrm{WO}(\Q)=\{S\in2^{\Q}:(S,<)\text{ is well-ordered}\}$ is $\PiC$-complete. By the first display of \cref{csf:thm:fixed-hard}, $E^{(1)}_{P_-}=2^{\Q}\setminus\mathrm{WO}(\Q)$, so the analytic-completeness of $E^{(1)}_{P_-}$ is the complement form of that lemma, transported by \cref{csf:thm:rank-one-classification}. Both arguments are the standard one and both are credited; the manuscript's reduction adds the fixed well-ordered prefix $W$, so that every target field has transcendence degree $\aleph_0$. That report's Theorem~18.4 (\lbl{pma:bor:thm:dichotomy}) also shows that $\mathrm{WO}(L)$ is Borel for scattered countable $L$ and $\PiC$-complete otherwise. Hence, as an observation of the write, if the rank-one parameters are restricted to subsets of a fixed countable $L\subseteq\Q$, the locus of $P_-$ is Borel exactly when $L$ is scattered; this bears on, but does not answer, Question~\ref{csf:q:fixedsource}.
\end{writenote}

\subsection{A second tree transform: copies of the rational order}

A related construction gives a self-contained hardness proof for the
predicate that a suborder of $\Q$ contains a copy of $\Q$. It will sharpen
the rare-ambient theorem in the larger family.

For a countable linear order $L$, let
\[
 Z(L)=\bigoplus_{l\in L}\Z
\]
with the reverse-lexicographic order: a nonzero vector is positive when
its coefficient at the largest nonzero rank is positive.

\begin{lemma}[An integer-sum transform]
\label[lemma]{csf:lem:integer-transform}
The order $Z(L)$ contains a copy of $\Q$ if and only if $L$ is not a
well-order.
\end{lemma}
\begin{proof}
If $L$ contains a decreasing sequence with range $D$ of type $\omega^*$,
consider the subgroup $Z(D)$. For $f<g$ in it, choose a rank $d\in D$
strictly below the largest nonzero rank of $g-f$. The unit vector $e_d$
then satisfies $0<e_d<g-f$, so $f<f+e_d<g$. This subgroup is a countable
dense linear order without endpoints, and hence isomorphic to $\Q$.

Conversely, suppose $u:\Q\emb Z(L)$. Choose rational endpoints $a<b$ and
let $l$ be the largest rank of $u(b)-u(a)$. Throughout the interval
$u(a)<u(r)<u(b)$, all coordinates above $l$ equal the corresponding
endpoint coordinates. The coordinate at $l$ can take only finitely many
integer values between its two endpoint values. Among the infinitely many
rational $r\in(a,b)$, choose two $a'<b'$ with the same coefficient at $l$.
Their nonzero difference has largest rank $l'<l$. Repeating this argument
inside $(a',b')$ produces an infinite strictly decreasing sequence of ranks
in $L$. Thus $L$ is not well ordered.
\end{proof}

Fix the countable ordered group
$G=Z((\N^{<\N},<_{\KB}))$, and choose once and for all an order embedding
$\psi:G\emb\Q$. For a tree $T$, define
\[
 R_T=\{\psi(g):g\in G,\ \supp(g)\subseteq T\}.
\]
Each membership coordinate of $R_T$ is either constantly false or a finite
conjunction of membership bits of $T$. Therefore $T\mapsto R_T$ is
continuous. The induced order on $R_T$ is exactly $Z((T,<_{\KB}))$.

\begin{theorem}[A second analytic-complete fixed source]
\label{csf:thm:rational-source}
The set
\[
 \mathcal N=\{S\subseteq\Q:(\Q,<)\emb(S,<)\}
           = E^{(1)}_{\Kone}
\]
is analytic-complete. Indeed,
\[
 T\in\IF\quad\Longleftrightarrow\quad R_T\in\mathcal N.
\]
\end{theorem}
\begin{proof}
The equivalence follows from \cref{csf:lem:kb,csf:lem:integer-transform}. The
reduction is continuous, and the target condition is analytic because an
order embedding is a countable function satisfying countably many tests.
The field equality follows from \cref{csf:thm:rank-one-classification}.
\end{proof}

In standard terminology, $\mathcal N$ consists of the non-scattered
suborders of $\Q$. The proof above establishes the required completeness
directly, rather than assuming an additional classification theorem about
scattered orders.

\subsection{Witness selection and its exact limitation}

Fix an enumeration $(a_n)$ of a source $P$. A function
$f:\N\to\Kstar$ codes an embedding $P\emb K_A$ if it is injective,
respects the field tables and constants of $P$, and all its values belong
to $K_A$. Each finite field-table test is clopen in the function space.
By finite supports, each test $f(n)\in K_A$ is also clopen. Their countable
intersection is a closed witness relation
\[
 \mathcal W_P\subseteq 2^I\times(\Kstar)^\N.
\]
Projection proves that every fixed-source locus is analytic. The analogous
relation for two varying fields is at least Borel, which suffices for the
analytic upper bound on embeddability.

\begin{corollary}[No total Borel solver]
\label[corollary]{csf:cor:no-selector}
There is no total Borel function on $2^{\Q}$ that outputs a valid embedding
$P_-\emb\Kone_S$ whenever one exists, with arbitrary outputs permitted
when none exists. There is likewise no Borel decision procedure for
membership in $E^{(1)}_{P_-}$.
\end{corollary}
\begin{proof}
For a total Borel candidate $f$, testing
$(S,f(S))\in\mathcal W_{P_-}$ is Borel. Under the asserted correctness
condition this test holds exactly on $E^{(1)}_{P_-}$, contradicting
\cref{csf:thm:fixed-hard}. The decision assertion is immediate.
\end{proof}

This statement is deliberately about \emph{total Borel} solvers. It neither
rules out non-Borel choice functions nor claims that a partial procedure
with a non-Borel domain is impossible.

\subsection{Why colors will be necessary}

Embeddability on $\Cone$ is Borel bireducible with embeddability of
uncolored countable linear orders, allowing finite domains. In one
direction use \cref{csf:thm:rank-one-classification}; in the other use the
rational placement above. The known better-quasi-order behavior of
countable linear orders prevents this from being a complete analytic
quasiorder \cite{laver,cmms}. Thus \cref{csf:thm:fixed-hard} is not a proof of
quasiorder universality. The larger family $\Cstar$ supplies countably
many intrinsically distinguishable Archimedean components, which removes
this obstruction.

\section{Quadratic colors and complete analytic embeddability}
\label{csf:sec:colors}

\subsection{An explicit antichain of ordered divisible groups}

For $n\in\N$, regard
\[
 H_n=\Q(\sqrt{p_n})=\Q+\Q\sqrt{p_n}
\]
as an ordered additive group. These groups are countable, divisible, and
Archimedean.

\begin{lemma}[Scalar rigidity]
\label[lemma]{csf:lem:hion}
If $H,H'\subseteq\R$ are additive subgroups with $1\in H$, and
$f:H\emb H'$ is an ordered additive embedding, then
$f(x)=\lambda x$ for every $x\in H$, where $\lambda=f(1)>0$.
For divisible subgroups the same proof is expressed with rational
multiples; more generally integer comparisons suffice.
\end{lemma}
\begin{proof}
For rational $r=m/n$ and $x\in H$, the inequality $r<x$ is equivalent to
$m<nx$. After applying $f$, it is equivalent to
$m\lambda<nf(x)$, or $r<f(x)/\lambda$. Hence $x$ and $f(x)/\lambda$
have the same rational lower cut in $\R$, so they are equal.
\end{proof}

This is the special case of Hion's lemma needed here; compare
\cite[Lemma 2.2]{cmms}.

\begin{lemma}[Quadratic-color antichain]
\label[lemma]{csf:lem:colors}
An ordered additive embedding $H_n\emb H_m$ exists if and only if $n=m$.
\end{lemma}
\begin{proof}
Suppose $\lambda H_n\subseteq H_m$ is such an embedding by
\cref{csf:lem:hion}. Since $\lambda=\lambda\cdot1$ belongs to the field $H_m$
and is nonzero, division by $\lambda$ gives $H_n\subseteq H_m$. Both
fields have degree two over $\Q$, so they are equal. Distinct prime
quadratic fields are unequal: writing
$\sqrt{p_n}=a+b\sqrt{p_m}$ with $a,b\in\Q$ and squaring forces
$a=0$ and $p_n/p_m=b^2$, which is impossible for distinct primes. The
converse uses the identity map.
\end{proof}

\subsection{Putting a colored order into the fixed field}

A countably colored linear order is a pair $L=(<_L,c_L)$ on $\N$, where
$c_L:\N\to\N$. Use the Borel rational placement $e_L$ of its uncolored
order, and define
\begin{equation}
 A_L=\{(e_L(i),1),(e_L(i),\sqrt{p_{c_L(i)}}):i\in\N\},
 \qquad M_L=K_{A_L}.
 \label{csf:eq:color-code}
\end{equation}
At rank $e_L(i)$ the value group of $M_L$ has Archimedean component
$H_{c_L(i)}$. In particular, all fields $M_L$ lie inside the same $\Kstar$;
the ambient field is not reconstructed for each input.

The map $L\mapsto A_L$ is Borel. For a given $(q,b)$, membership is a
countable union of conditions involving $e_L(i)=q$ and one color value.
Thus $L\mapsto M_L$ is Borel by \cref{csf:thm:cantor}.

\begin{theorem}[Exact colored-order coding]
\label{csf:thm:colored-coding}
For countably colored countable linear orders $L,L'$,
\begin{align*}
 M_L\emb M_{L'}
 &\quad\Longleftrightarrow\quad
 L\emb L'\text{ by an order- and color-preserving injection},\\
 M_L\cong M_{L'}
 &\quad\Longleftrightarrow\quad
 L\cong L'\text{ as colored orders}.
\end{align*}
No color predicates are added to the field language.
\end{theorem}
\begin{proof}
Let $f:M_L\emb M_{L'}$. Its value-group map induces an order embedding
of rank chains, and hence an order injection $h:L\emb L'$. At each rank,
\cref{csf:lem:valuation-functor} gives an ordered group embedding
\[
 H_{c_L(i)}\emb H_{c_{L'}(h(i))}.
\]
By \cref{csf:lem:colors}, the colors are equal. An isomorphism of fields makes
the rank map an isomorphism.

Conversely, for a color-preserving order injection $h$, map the two
independent generators at $e_L(i)$ to the corresponding two generators
at $e_{L'}(h(i))$. For a finite monomial exponent, this relabels the rank
coordinates while preserving the real coefficient within each coordinate.
The order of exponents and the leading sign of every polynomial are
therefore preserved. Algebraic independence gives an ordered embedding of
the rational-function fields; real-closure uniqueness extends it to the
required field embedding. A bijective $h$ gives an isomorphism.
\end{proof}

\begin{theorem}[Complete analytic quasiorder in a closed subfield family]
\label{csf:thm:universality}
Embeddability on $\Cstar$ is a complete analytic quasiorder. Its associated
bi-embeddability relation is a complete analytic equivalence relation.
The hardness already holds on the Borel-coded fields $M_L$, all of which
have transcendence degree $\aleph_0$ and residue field $k$.
\end{theorem}
\begin{proof}
Embeddability is analytic by coding a map between countable fields and
imposing the countable collection of field-table conditions. By
\cref{csf:thm:colored-coding}, the Borel map $L\mapsto M_L$ reduces colored-order
embeddability to it. Louveau's completeness theorem for the former
quasiorder \cite{marcone,cmms} proves the desired universality.

For bi-embeddability, let $E$ be any analytic equivalence relation. As an
analytic quasiorder it reduces to embeddability via a Borel map $f$.
Symmetry of $E$ then gives
\[
 xEy\quad\Longleftrightarrow\quad f(x)\emb f(y)\ \&\ f(y)\emb f(x).
\]
The converse follows from either embedding and the reduction property.
Bi-embeddability is analytic, being the intersection of the two analytic
embedding relations. Finally, each infinite-domain code
\eqref{csf:eq:color-code} has countably infinitely many independent
coordinates, so \eqref{csf:eq:trdeg} and \cref{csf:thm:valuation} give the stated
invariants.
\end{proof}

Together with \cref{csf:cor:isomorphism,csf:thm:cantor}, this proves
\cref{csf:thm:main}. The novelty claim is limited: the broad universality theorem
was known; the present proof places an explicit universal coding in a
compact coordinate family of actual elementary subfields of a single
specified countable surreal field.

\begin{writenote}[Added 4 October 2026, batch~90]
The write checked the cited statements of \cite{cmms} (Lemma~2.2, Theorem~7.1, Proposition~7.2, Corollary~7.3; PDF pages 33--35 of version~2); they say what is attributed to them here. See \cref{csf:sec:provenance}.
\end{writenote}

\section{The finite-transcendence-degree threshold}
\label{csf:sec:threshold}

The finite-support argument imposes a sharp restriction on fixed-source
existence questions even when the global embeddability quasiorder is
maximally complicated.

\begin{theorem}[Finite sources have finite-coordinate witnesses]
\label{csf:thm:finite-source}
Let $P$ be a real closed field of finite transcendence degree over its
canonical copy of $k$. Then
\begin{equation}
 E^\star_P=\bigcup_{\substack{U\subseteq I\text{ finite}\\P\emb K_U}}
             \{A:U\subseteq A\}.
 \label{csf:eq:open-locus}
\end{equation}
In particular $E^\star_P$ is an open upward-closed subset of $2^I$.
The same formula holds for the rank-one family with finite $U\subseteq\Q$.
\end{theorem}
\begin{proof}
Suppose $f:P\emb K_A$. Choose a finite transcendence basis
$a_1,\ldots,a_d$ for $P/k$. Let
\[
 U=D(f(a_1))\cup\cdots\cup D(f(a_d)).
\]
Then $U$ is finite and contained in $A$. All $f(a_j)$ lie in $K_U$.
Every element of $f(P)$ is algebraic over $k(f(a_1),\ldots,f(a_d))$, and
hence over $K_U$. Since $K_U$ is real closed, it has no proper algebraic
ordered extension inside $\Kstar$. Therefore $f(P)\subseteq K_U$.

Conversely, any embedding into $K_U$ composes with inclusion into $K_A$
whenever $U\subseteq A$. Each upper cone on the right side of
\eqref{csf:eq:open-locus} is clopen, so the union is open. The rank-one argument
is identical. When $d=0$, the source is $k$ and the witness is $U=\varnothing$.
\end{proof}

\begin{proposition}[Infinite sources have no interior]
\label[proposition]{csf:prop:infinite-source}
If $\trdeg_k P=\aleph_0$, both $E^\star_P$ and $E^{(1)}_P$ have empty
interior. Thus, in either family, a nonempty fixed-source locus is open if
and only if its source has finite transcendence degree.
\end{proposition}
\begin{proof}
Finite subsets of the coordinate set are dense in its power set: every
basic open cylinder only prescribes finitely many present and absent
coordinates, and the required present coordinates alone give a finite
member of the cylinder. By \eqref{csf:eq:trdeg}, no field from a finite
coordinate set can contain an embedded copy of $P$. Every nonempty open
set therefore meets the complement of the locus. Combine this with
\cref{csf:thm:finite-source}.
\end{proof}

The nonempty qualification matters. A field not embeddable in the chosen
ambient field has empty, hence open, locus regardless of its transcendence
degree. We do not assert that $\Kstar$ contains every countable real closed
field; the fixed residue field alone precludes such a claim.

\begin{example}[The simplest finite cases]
If $U\subseteq\Q$ has $n<\omega$ elements, then
\[
 E^{(1)}_{\Kone_U}=\{S:|S|\ge n\}.
\]
Every $n$-element linear order is a chain, and it embeds into $S$ exactly
when $S$ has at least $n$ elements. For $n\ge1$, the complement is the
closed set of subsets of size at most $n-1$.
\end{example}

\begin{remark}[Open is not an effective claim]
\eqref{csf:eq:open-locus} does not furnish an algorithm deciding which finite
$U$ admit an embedding of an arbitrarily specified $P$. It proves a
boldface topological fact: every positive instance has some finite
coordinate neighborhood consisting entirely of positive instances. No
uniform computability of the displayed union is claimed.
\end{remark}

\section{Generic collapse in the rank-one family}
\label{csf:sec:generic-one}

Equip $2^{\Q}$ with the product probability measure in which each rational
is included independently with probability $1/2$. This choice is a measure
on generator parameters, not on the surreal field.

\begin{theorem}[A comeager conull isomorphism type]
\label{csf:thm:generic-collapse}
The set
\[
 \mathcal G=\{S\subseteq\Q:S\cap(a,b)\ne\varnothing
                 \text{ for every }a<b\text{ in }\Q\}
\]
is a dense $G_\delta$ of measure one. For every $S\in\mathcal G$,
$\Kone_S\cong\Kone$. Consequently every nonempty $E^{(1)}_P$ contains
$\mathcal G$ and is comeager and conull.
\end{theorem}
\begin{proof}
For fixed rational $a<b$, the condition of meeting $(a,b)$ is open and
dense: any finite cylinder can be extended by including a still
unmentioned rational in that interval. It has probability one because
$(a,b)\cap\Q$ is infinite, and the probability of missing its first $m$
points is $2^{-m}$. There are only countably many rational intervals, so
their intersection is a dense $G_\delta$ of measure one.

Every $S\in\mathcal G$ is a countable dense linear order without endpoints.
Cantor's back-and-forth theorem gives $S\cong\Q$. Now apply
\cref{csf:thm:rank-one-classification}.

If $E^{(1)}_P$ is nonempty, then $P$ embeds into some $\Kone_T\subseteq\Kone$.
For $S\in\mathcal G$, composing with an isomorphism
$\Kone\cong\Kone_S$ proves $S\in E^{(1)}_P$.
\end{proof}

The last assertion also shows density, but there is a useful elementary
version. A basic cylinder forbids only finitely many rationals, say $V$.
The set $\Q\setminus V$ has order type $\Q$, meets that cylinder after
respecting its required coordinates, and yields another copy of the
ambient field.

\begin{corollary}[Hardness concentrated on an exceptional set]
\label[corollary]{csf:cor:hard-generic}
The analytic-complete set $E^{(1)}_{P_-}$ is comeager and conull. Its
coanalytic-complete complement is meager and null. In particular, maximum
analytic decision complexity can coexist with probability-one agreement
on the answer.
\end{corollary}

There is no conflict with \cref{csf:prop:infinite-source}: a comeager set may
have empty interior. Indeed, every finite coordinate pattern can still be
completed to a finite-rank field that does not admit $P_-$.

\section{Two generic regimes in the universal family}
\label{csf:sec:generic-star}

The larger family $\Cstar$ does not inherit the same generic collapse.
Its additional coefficient coordinates can remove an entire Archimedean
component type while leaving almost every rank active. This produces two
opposite fixed-source phenomena within one compact space.

\subsection{Dense loci and general zero--one laws}

Use independent fair coin choices on the countable coordinate set
$I=\Q\times B$. For $r\in\Q$, translation of the first coordinate gives
\[
 \tau_r(A)=\{(q+r,b):(q,b)\in A\}.
\]
The monomial relabeling $x_{q,b}\mapsto x_{q+r,b}$ extends to an isomorphism
$K_A\cong K_{\tau_r(A)}$: it preserves the finite reverse-lexicographic
exponent order. Thus each fixed-source locus $E^\star_P$ is invariant under
all $\tau_r$.

\begin{theorem}[Density and zero--one laws]
\label{csf:thm:zero-one}
For every countable real closed $P$, $E^\star_P$ is analytic. If it is
nonempty, it is dense. It is either meager or comeager, and its product
measure is either zero or one. The same zero--one laws hold for every
translation-invariant analytic subset of $2^I$.
\end{theorem}
\begin{proof}
Analyticity follows from the closed witness relation in
\cref{csf:sec:descriptive}. Suppose $P\emb\Kstar$, as holds whenever the locus
is nonempty. In a basic cylinder let $V\subset I$ be the finite set of
forbidden coordinates, and set $A=I\setminus V$. All columns at ranks
outside the finite projection of $V$ onto $\Q$ are complete. That cofinite
rank set has order type $\Q$, so rank relabeling embeds $\Kstar$ into
$K_A$. Hence $P\emb K_A$, and $A$ lies in the cylinder. This proves density.

We prove the category statement for a translation-invariant set $E$ with
the Baire property, which every analytic set has. If $E$ is nonmeager, it
is comeager on some nonempty basic cylinder $U$. Given any other nonempty
basic cylinder $V$, choose $r\in\Q$ so that the finitely many rank
coordinates mentioned by $\tau_r(U)$ are disjoint from those mentioned
by $V$. Then $\tau_r(U)\cap V$ is nonempty and open. Invariance makes $E$
comeager on $\tau_r(U)$. If the complement of $E$ were nonmeager, it would
be comeager on some such $V$, contradicting the Baire theorem on their
intersection. Therefore $E$ is comeager.

For measure, analytic sets are measurable in the completion of the product
measure. Approximate $E$ in measure by a finite-coordinate event $C$,
with $\mu(E\mathbin\triangle C)<\varepsilon$. Choose a rational translation
$r$ making the rank coordinates of $C$ and $\tau_r(C)$ disjoint. These two
events are independent, so
\[
 \mu(C\cap\tau_r(C))=\mu(C)^2.
\]
Since $E=\tau_r(E)$, the symmetric difference between $E$ and
$C\cap\tau_r(C)$ has measure at most $2\varepsilon$. Letting
$\varepsilon\to0$ gives $\mu(E)=\mu(E)^2$, and hence $\mu(E)\in\{0,1\}$.
\end{proof}

The category and measure conclusions are separate dichotomies. The theorem
does not say they agree for every source. The two explicit sources below
do have matching category and measure behavior.

\subsection{The reverse-\texorpdfstring{$\omega$}{omega} source remains generically present}

For $A\subset I$, let
\[
 J_A=\{q\in\Q:A_q\ne\varnothing\},\qquad
 A_q=\{b\in B:(q,b)\in A\}.
\]

\begin{proposition}
\label[proposition]{csf:prop:rank-one-in-star}
For $S\subseteq\Q$ and $A\subseteq I$,
\[
 \Kone_S\emb K_A\quad\Longleftrightarrow\quad(S,<)\emb(J_A,<).
\]
Consequently $E^\star_{P_-}$ is analytic-complete, comeager, and conull.
\end{proposition}
\begin{proof}
Necessity follows from the intrinsic rank map. For sufficiency, take an
order injection $h:S\emb J_A$ and choose $b_q\in A_{h(q)}$ for each
$q\in S$. Send $t_q$ to $x_{h(q),b_q}$. Multiplication by the positive real
$b_q$ preserves the coefficient sign at the largest differing rank, so
this map preserves signs of polynomials in the source generators.
Algebraic independence and real-closure extension finish the embedding.

The face $S\mapsto S\times\{1\}$ continuously reduces
$E^{(1)}_{P_-}$ to $E^\star_{P_-}$, proving analytic-completeness. For each
fixed $q$, the condition $A_q\ne\varnothing$ is open dense and has
probability one, because $B$ is infinite. The intersection over $q\in\Q$
is comeager and conull, and on it $J_A=\Q$. Every rank-one source embeds
there, in particular $P_-$.
\end{proof}

\subsection{The whole ambient field is generically absent}

Put $H=\operatorname{span}_{\Q}B$. This is the Archimedean component of
$\Gamma_I$ at every rank.

\begin{lemma}[A complete component cannot fit into a partial column]
\label[lemma]{csf:lem:full-component}
For $C\subseteq B$, there is an ordered additive embedding
\[
 H\emb\operatorname{span}_{\Q}C
\]
if and only if $C=B$.
\end{lemma}
\begin{proof}
Let $f$ be an embedding. By \cref{csf:lem:hion}, $f(x)=\lambda x$
with $\lambda=f(1)>0$. Write
\[
 \lambda=a_0+\sum_{i\in F}a_i\sqrt{p_i},\qquad a_i\in\Q,
\]
with $F$ finite. Choose $j\notin F$. Since $\sqrt{p_j}\in H$,
\[
 \lambda\sqrt{p_j}
   =a_0\sqrt{p_j}+\sum_{i\in F}a_i\sqrt{p_ip_j}
 \in\operatorname{span}_{\Q}C\subseteq H.
\]
Distinct squarefree radical monomials are independent by
\cref{csf:lem:radicals}. In particular, no nonzero coefficient of
$\sqrt{p_ip_j}$ can occur in an element of $H$. Thus all $a_i$ for
$i\in F$ vanish, and $\lambda=a_0$ is a nonzero rational.
Now $\lambda\cdot1\in\operatorname{span}_{\Q}C$ forces $1\in C$;
and $\lambda\sqrt{p_n}\in\operatorname{span}_{\Q}C$ forces
$\sqrt{p_n}\in C$ for every $n$. Therefore $C=B$. The converse is the
identity embedding.
\end{proof}

Define the complete-column rank set
\[
 D_A=\{q\in\Q:A_q=B\}.
\]

\begin{theorem}[Exact whole-ambient locus]
\label{csf:thm:rare-ambient}
For $A\subseteq I$,
\begin{equation}
 \Kstar\emb K_A\quad\Longleftrightarrow\quad(\Q,<)\emb(D_A,<).
 \label{csf:eq:whole-locus}
\end{equation}
The locus $E^\star_{\Kstar}$ is analytic-complete and dense, but meager
and of measure zero. Its complement is coanalytic-complete.
\end{theorem}
\begin{proof}
If $\Kstar\emb K_A$, the rank map sends each source component $H$ into
the component at its image rank. By \cref{csf:lem:full-component}, that target
column must be complete. Thus the rank map embeds $\Q$ into $D_A$.
Conversely, an order injection $h:\Q\emb D_A$ permits the relabeling
$x_{q,b}\mapsto x_{h(q),b}$ for every $b\in B$. This extends to an embedding
of real closures and proves \eqref{csf:eq:whole-locus}.

The locus is nonempty because $\Kstar$ embeds into itself, and it is dense
by \cref{csf:thm:zero-one}. For each $q$, the set
$\{A:A_q=B\}$ is closed and nowhere dense: a finite cylinder can always
be refined by omitting one as-yet unmentioned coordinate in that column.
Its measure is zero, since the probability of including its first $m$
coordinates is $2^{-m}$. Therefore
\[
 \{A:D_A\ne\varnothing\}
   =\bigcup_{q\in\Q}\{A:A_q=B\}
\]
is meager and null. The locus in \eqref{csf:eq:whole-locus} is a subset of this
set, so it too is meager and null.

Finally, the map $S\mapsto S\times B$ is continuous and has
$D_{S\times B}=S$. By \eqref{csf:eq:whole-locus} it reduces the
analytic-complete set $\mathcal N$ of \cref{csf:thm:rational-source} to
$E^\star_{\Kstar}$. Analyticity was already proved in
\cref{csf:thm:zero-one}. This establishes analytic-completeness and the
corresponding completeness of the complement.
\end{proof}

\begin{corollary}[Two infinite sources, opposite answers almost surely]
The infinite-transcendence-degree sources $P_-$ and $\Kstar$ both have
dense embedding loci in $\Cstar$. The first locus is comeager and conull;
the second is meager and null. Both are analytic-complete and both have
empty interior.
\end{corollary}

This proves the remaining clauses of \cref{csf:thm:main-source}. It also shows
that the finite-versus-infinite threshold controls \emph{openness}, not
probability. Infinite sources alone already exhibit both generic regimes.

\section{Examples, diagnostic counterexamples, and interpretation}
\label{csf:sec:examples}

\subsection{A two-scale calculation}

For $q<r$, set $u=\omega^{\omega^q}$ and $v=\omega^{\omega^r}$. For every
positive integer $n$,
\[
 u^n=\omega^{n\omega^q}<\omega^{\omega^r}=v.
\]
This is the field-level trace of the strict order between two
Archimedean ranks in the value group. For a polynomial such as
\[
 P(u,v)=3u^7v^2-5u^{100}v+2u^3-9,
\]
the $v^2$ term has the largest exponent regardless of the larger $u$
exponent in the next term. Hence $P(u,v)>0$. Replacing $(q,r)$ by any
increasing pair preserves the same comparison. This finite dominance is
the mechanism behind the uncolored embeddings.

\subsection{A color is a ratio type, not a named constant}

At one rank $q$, the pair
\[
 \omega^{\omega^q},\qquad \omega^{\sqrt2\,\omega^q}
\]
generates, after real closure, a component $\Q+\Q\sqrt2$. Replacing
$\sqrt2$ by $\sqrt3$ gives a different ordered component. Although both
are abstractly two-dimensional rational vector spaces, no order-preserving
additive embedding exists between them. The source and target field
languages contain neither a color predicate nor a distinguished value-group
basis. The distinction survives because their ordered ratio structures are
intrinsic to the fields' natural valuations.

\subsection{Why finite rank cannot secretly code the fixed hard source}

A finite coordinate set produces finite transcendence degree. It cannot
contain an embedded copy of $P_-$, even if its generators have very large
surreal birthdays or complicated real coefficients. Enlarging the
complexity of finitely many terms is not the same as supplying infinitely
many independent rank levels. This guards against an erroneous attempt
to put the tree-hard fixed-source problem inside a finitely generated
real closure.

\subsection{What this does and does not say about algorithms}

An analytic-completeness theorem is a non-Borel definability obstruction.
It is stronger than failure of one proposed normal form, but it is not a
numerical runtime lower bound and does not by itself prove a particular
finite-input problem undecidable. Effective consequences require an
explicit computable coding of field presentations and of the reductions.
The present paper supplies Borel and continuous codings, not a full
computable presentation of the surreal ambient field.

Conversely, the clopen membership tests in \cref{csf:thm:cantor} show that the
coding map itself is not the source of the descriptive difficulty.
Individual elements require finitely many coordinates. It is the
existential quantification over an entire embedding, and the infinite
order structure that embedding must preserve, which creates the
non-Borel locus.

\section{Implications for mathematical evaluation and ProveIt}
\label{csf:sec:formalization}

\subsection{A sampling warning with a proof}

The link to evaluation is a mathematical deduction, not an empirical
claim about any particular benchmark. In the rank-one family, sampling
each generator independently almost surely produces a field isomorphic
to $\Kone$. A test that merely answers ``yes'' to the fixed-source
question for any $P\emb\Kone$ succeeds with probability one, even when
its positive locus is analytic-complete. Such a test has not learned the
exceptional structural cases.

In the full family, the reverse constant answer is almost surely correct
for the source $\Kstar$: its embedding locus is null. Neither behavior is
a satisfactory test of the exact equivalence in
\cref{csf:thm:rank-one-classification,csf:thm:colored-coding,csf:thm:rare-ambient}.
A structure-sensitive evaluation would deliberately include well-founded
versus ill-founded tree codes, incompatible quadratic components,
finite-transcendence-degree obstructions, and embeddings with nontrivial
rank maps. Finite tree truncations alone do not certify an infinite
branch, so they must not be mislabeled as tests of the full analytic
existence theorem.

\subsection{Repository snapshot and reusable boundaries}

The repository inspected for this article was pinned to
\begin{center}
\texttt{7c0f2d9f92c3d51ec85703bed1022924a4b7359b}.
\end{center}
The inspected surreal README distinguishes formal declarations from
AI-assisted unrefereed reports \cite{proveit}. Search and source inspection
identified the following concrete integration points:
\begin{center}
\begin{tabular}{P{0.53\textwidth}P{0.39\textwidth}}
\toprule
Repository path, relative to \texttt{Algebra/SurrealNumbers/} & Relevant boundary\\
\midrule
\texttt{Surreal/Foundations/}\newline
\texttt{SignSequenceRealClosed.lean} & Real closedness of the actual
sign-sequence surreal field.\\[4pt]
\texttt{Surreal/HahnSeries/RealClosed.lean} & Set-sized ordered Hahn-field
real closedness.\\[4pt]
\texttt{SurrealAudit.lean} & Existing axiom-audit infrastructure.\\[4pt]
\texttt{docs/surreal/}\newline
\texttt{surreal-self-embeddings/README.md} & Research background on rank,
Hahn lifts, and size-safe foundations; not a formal proof of this article.\\
\bottomrule
\end{tabular}
\end{center}
The declaration name \texttt{signSequenceIsRealClosed} occurs in the
inspected foundation source. No claim is made that a build of the
repository, or a kernel check of any new theorem here, was performed.

\begin{writenote}[Added 4 October 2026, batch~90]
At the write the three Lean files of the table are byte-identical to the pin, and the descriptions are accurate (\cref{csf:sec:provenance}). The table's paths are relative to \repo{Algebra/SurrealNumbers/}. Beyond them, the only formal statement bearing on this report is the multiquadratic independence behind \cref{csf:lem:radicals} (note after that lemma); everything else here is unformalized, and the plan below remains a proposal.
\end{writenote}

\subsection{A staged formalization plan}

The most reusable theorem is \cref{csf:thm:support}. It should be formalized
first as a field-theoretic statement independent of surreal numbers.
One needs an indexed algebraically independent family, a relative
algebraicity predicate, finite character, and exchange. Its conclusion
then gives continuity into the product of Boolean coordinates, arbitrary
intersection formulas, and directed-union formulas.

A second stage treats finite-support ordered groups. For a linear order
$L$ and ordered Archimedean components $H_l$, use finitely supported
functions with reverse-lexicographic order. Prove the rank computation and
the induced maps on convex quotients. The quadratic-color obstruction is
a finite algebraic subproject, separate from the higher descriptive
complexity inputs.

A third stage constructs the coordinate monomial field and its relative
real closure in a set-sized ordered real closed ambient field. The exact
value group, residue field, and generator-relabeling maps provide the
bridge to \cref{csf:thm:colored-coding}. The surreal normal-form interface is
needed only to instantiate the independent ordered monomial family.

Finally, implement the standard Borel spaces and the reductions. The
analytic-completeness theorem for colored orders is a substantial external
library dependency; it should either be proved, imported from an audited
development, or left as an explicit hypothesis in a clearly named
conditional theorem. Hiding it as an unexplained axiom would defeat the
point of the proof boundary.

\subsection{Finite verification included with this article}

The accompanying Python program performs exact finite checks of
reverse-lexicographic dominance, polynomial multiplication signs under
increasing coordinate embeddings, finite-order embedding tests, quadratic
radical character orthogonality, and the finite rational-placement
algorithm. It uses rational and integer arithmetic, not floating-point
surreal approximations. Its output records the seed and actual check
counts.

These tests check implementation conventions and expose sign-orientation
mistakes. They do not prove algebraic independence for infinitely many
generators, the finite-support theorem, analytic completeness, or the
category and measure theorems. Those are supported by the mathematical
proofs and the stated external inputs.

\section{Further research questions}
\label{csf:sec:questions}

The following questions are proposals for continued work. Their worldwide
open status is not asserted. They are organized around gaps not settled
by the present proofs, rather than restating the theorems already obtained.

\begin{question}[Exact fixed-source complexity]\label[question]{csf:q:fixedsource}
For $P\emb\Kone$, classify the Borel or projective complexity of
$E^{(1)}_P$. For sources of the form $\Kone_S$, this becomes the complexity
of $\{T:S\emb T\}$. Determine which infinite source orders give Borel
loci and which give analytic-complete loci. The present finite/infinite
threshold decides openness, not this finer classification.
\end{question}

\begin{question}[Generic isomorphism in the larger family]\label[question]{csf:q:genericiso}
Determine the generic behavior of isomorphism on $\Cstar$. The ambient
isomorphism class is meager and null by \cref{csf:thm:rare-ambient}, so the
rank-one collapse does not identify the generic type here. Is every
isomorphism class meager or null? Which invariants describe independent
random component subspaces $H_A(q)$?
\end{question}

\begin{question}[Generic-presence criterion]\label[question]{csf:q:presence}
Give an intrinsic criterion on a countable source $P\emb\Kstar$ deciding
whether $E^\star_P$ is comeager or meager. Is there a corresponding criterion
in terms of the source's ordered value-group components and their finite
supports? Do the category and measure alternatives always agree for these
particular fixed-source loci?
\end{question}

\begin{question}[Effective realization]\label[question]{csf:q:effective}
Construct a uniform computable presentation of the rational monomial
fields and their ordered real closures for suitable computable rank orders.
Then determine the lightface complexity of the fixed-source problem.
This requires controlling effective real-closure construction; the Borel
proof alone is insufficient.
\end{question}

\begin{question}[Canonical support computation]\label[question]{csf:q:support}
For an algebraic element specified by an isolating polynomial over the
coordinate rational-function field, develop an algorithm computing its
minimal support $D(a)$. Which operations permit efficient support
cancellation tests? This is directly relevant to symbolic implementations
and to a formal proof of the continuity theorem.
\end{question}

\begin{question}[Quantitative birthday bounds]\label[question]{csf:q:birthday}
Find an explicit ordinal $\beta$ containing the displayed $\Kstar$, ideally
with a proof stable under real closure. Compare the birthday cost of this
double-monomial realization with other realizations of the same ordered
value groups. The existential set bound in \cref{csf:sec:foundations} is not
an optimality result.
\end{question}

\begin{question}[How much compact localization is general?]\label[question]{csf:q:localization}
For other model-complete theories with an algebraic closure pregeometry,
identify sufficient conditions for a closed coordinate family of elementary
submodels to support complete analytic embeddability. The finite-support
lemma is general, but real closure and monomial order preservation supply
additional structure that cannot simply be omitted.
\end{question}

\begin{question}[Individual embeddings and automorphism groups]\label[question]{csf:q:embeddings}
For a fixed rank-one or colored field, classify all embeddings inducing a
prescribed rank map. Determine the kernel of the map from its automorphism
group to automorphisms of the ordered component data, with the
pointwise-convergence topology on the countable domain.
\end{question}

\begin{question}[Additional named structure]\label[question]{csf:q:named}
What survives after naming selected generators, a monomial section, an
integer part, or an appropriate restricted exponential operation? Naming
all generators trivializes many maps, while exponential closure can destroy
the independent-coordinate construction. A useful result would identify
a nontrivial boundary between these extremes.
\end{question}

\begin{question}[Invariant universality versus completeness]\label[question]{csf:q:invariant}
Does the embeddability quasiorder restricted to a suitably chosen Borel
subfamily satisfy a stronger invariant-universality property compatible
with isomorphism? \cref{csf:thm:universality} proves completeness only; the
additional invariance requirements need a separate argument.
\end{question}

\begin{question}[Changing the sampling law]\label[question]{csf:q:sampling}
Replace the fair coins by coordinate-dependent inclusion probabilities.
Characterize when rank-one generic collapse persists, and when complete
columns in the larger family appear with positive probability. Infinite
products and interval-coverage conditions suggest explicit thresholds,
but a full source-dependent theorem remains to be developed.
\end{question}

\begin{question}[Proof-certified evaluation instances]\label[question]{csf:q:certified}
Build a library of finite, checkable implications drawn from the source
threshold, color incompatibility, and rank embeddings. Separate certificates
for finite algebraic facts from assumptions about infinite trees, and
measure whether a reasoning system respects that distinction. The goal is
not to pretend that finite samples decide an analytic-complete set.
\end{question}

\begin{writenote}[Added 4 October 2026, batch~90]
The write found no report of the collection that answers any of these twelve questions. Related material: for \cref{csf:q:fixedsource}, the note after \cref{csf:thm:fixed-hard} settles the source $P_-$ on parameter families restricted to subsets of a fixed countable order, by \lbl{pma:bor:thm:dichotomy}; for \cref{csf:q:effective}, \repo{foundations-and-computation/computable-surreals} gives effective countable real closed subfields of $\No$, but with value group $\Q$, not the families of ranks needed here; for \cref{csf:q:birthday}, the collection's birthday reports (\repo{foundations-and-computation/birthday-cutoffs-and-hereditary-sets}, \repo{surreal/gonshor-product-birthdays}, \repo{surreal/gonshor-laurent-birthdays}) are the places to look, but whether their bounds yield one for this~$\beta$ was not examined; for \cref{csf:q:embeddings}, \repo{surreal/surreal-self-embeddings} studies embeddings of all of $\No$, not of these countable fields. The labels \texttt{csf:q:*} were added at the write.
\end{writenote}

\section{Conclusion and claim status}

A countable surreal field can contain a closed Cantor family whose
membership coordinates are all clopen and whose embeddability quasiorder
is analytically universal. The two facts coexist because finite support
controls individual elements, whereas an embedding must preserve an
infinite hierarchy of valuation ranks and components.

The rank-one subfamily gives an exact dictionary between field embeddings
and order embeddings. It yields a fixed-source analytic-complete locus,
an open-locus theorem for finite transcendence degree, and a probability-one
isomorphism type. Adding quadratic coefficient colors both upgrades the
quasiorder complexity and creates a contrasting source that is almost
surely absent. The category and measure zero--one laws make this contrast
structural rather than anecdotal.

The claims established here are proof-level refinements of known
anti-classification machinery, not a claim to have first discovered the
unrestricted complexity of real closed fields. The manuscript is an
AI-assisted, unrefereed research draft. Its local theorems have detailed
proofs; independent mathematical review and any future formalization remain
separate validation steps.

\appendix
\section{Proof-dependency ledger}
\label{csf:app:ledger}

\begin{longtable}{P{0.25\textwidth}P{0.38\textwidth}P{0.28\textwidth}}
\toprule
Result & Mathematical dependencies & Status in this package\\
\midrule
\endfirsthead
\toprule
Result & Mathematical dependencies & Status in this package\\
\midrule
\endhead
Double-monomial independence & Finite normal-form comparison; multiquadratic
independence & Full proof, \cref{csf:lem:independence}.\\[4pt]
Exact values and residues & Natural valuation; minimum-term cancellation;
real closedness of $k$ & Full proof, \cref{csf:thm:valuation}.\\[4pt]
Finite supports & Finite character and field exchange & Full proof,
\cref{csf:thm:support}; no surreal-specific hypothesis.\\[4pt]
Cantor localization & Finite supports; compact-to-Hausdorff argument & Full
proof, \cref{csf:thm:cantor}.\\[4pt]
Rank-one classification & Intrinsic value-group rank; real-closure extension
& Full proof, \cref{csf:thm:rank-one-classification}.\\[4pt]
Fixed-source hardness & Rank classification; Kleene--Brouwer lemma;
completeness of $\IF$ & New reduction proved in full; $\IF$ completeness
credited.\\[4pt]
Colored-order coding & Scalar rigidity; quadratic antichain; intrinsic
components & Full proof, \cref{csf:thm:colored-coding}.\\[4pt]
Complete analytic quasiorder & Colored-order coding; Louveau's theorem &
Coding proved; external completeness input explicit.\\[4pt]
Finite-source threshold & Finite supports; algebraicity of source over a
finite basis & Full proof, \cref{csf:thm:finite-source,csf:prop:infinite-source}.\\[4pt]
Generic collapse & Rational interval coverage; countable dense-order
back-and-forth & Full proof, \cref{csf:thm:generic-collapse}.\\[4pt]
Zero--one laws & Translation symmetry; Baire property; finite-cylinder
measure approximation & Full proof, \cref{csf:thm:zero-one}.\\[4pt]
Rare whole-ambient embeddings & Radical independence; component map;
integer-sum tree transform; countable null union & Full proof,
\cref{csf:thm:rare-ambient}.\\[4pt]
Lean verification & Would require newly implemented statements and a
kernel run & Not supplied or claimed.\\
\bottomrule
\end{longtable}

\section{Finite checks and reproducibility}
\label{csf:app:checks}

The archive contains \texttt{article.tex}, \texttt{article.pdf}, a short
README, a source-and-claim audit, and the following small reproducibility
assets:
\begin{center}
\begin{tabular}{ll}
\toprule
\texttt{code/finite\_checks.py} & Standard-library exact finite checks.\\
\texttt{data/finite\_checks.json} & Recorded output of the delivered run.\\
\texttt{build.sh} & Repeats the checks and compiles the article.\\
\texttt{SHA256SUMS.txt} & Checksums of the delivered files.\\
\bottomrule
\end{tabular}
\end{center}
Run from the package directory:
\begin{lstlisting}
python3 code/finite_checks.py --output data/finite_checks.json
pdflatex -interaction=nonstopmode -halt-on-error article.tex
pdflatex -interaction=nonstopmode -halt-on-error article.tex
pdflatex -interaction=nonstopmode -halt-on-error article.tex
\end{lstlisting}
No font files, external data downloads, or nonstandard Python packages are
required. A normal TeX Live installation with the packages named in the
preamble is sufficient.

The finite polynomial representation stores exponent vectors in increasing
rank-coordinate order, so the largest index is the most significant one.
After combining equal monomials, the leading coefficient determines the
sign. Products are formed by exact convolution, including cancellation,
and checked against the ordered-field sign law and monotone reindexing.
The rational-placement tests enumerate finite orders and use exact
fractions. Radical checks implement independent sign characters, rather
than numerical approximations to square roots.

The recorded checks are deliberately modest evidence. A passed finite test
cannot certify the non-Borel conclusions, and no such interpretation of
\texttt{finite\_checks.json} should be made.

\begin{writenote}[Added 4 October 2026, batch~90]
In ProveIt the package is laid out differently. The archive's \texttt{article.pdf}, README and \texttt{SHA256SUMS.txt} are not shipped (the checksums were verified at placement; the PDF is replaced by a build of this text, the README by a report README); the source audit is \texttt{SOURCE\_AND\_CLAIM\_AUDIT.md}; \texttt{build.sh} is shipped as \texttt{code/build.sh}, where its change to its own directory makes it look for \texttt{code/code/finite\_checks.py} and \texttt{code/article.tex}, so it must be run on a copy in the delivered layout. Run in place with \texttt{--output data/finite\_checks.json}, the checks would overwrite the recorded file; write to a scratch path instead. At the write a rerun on a copy (Python~3.14.4) passed all 25,378 checks and reproduced \texttt{data/finite\_checks.json} except for the \texttt{python\_version} line (the recorded run used 3.13.5); on Windows the output has CRLF line ends. The report README gives the commands, and the archive can be recovered with \texttt{git show a162e4386:docs/incoming/Surreal\_Subfield\_Cantor\_Research.zip}.
\end{writenote}

\begin{thebibliography}{99}
\raggedright

\bibitem{glazer}
Elliot Glazer.
\emph{A Topological Tennenbaum Theorem}.
Preprint, 2023. \href{https://arxiv.org/abs/2311.13699}{arXiv:2311.13699}.

\bibitem{frontier}
Elliot Glazer et al.
\emph{FrontierMath: A Benchmark for Evaluating Advanced Mathematical
Reasoning in AI}.
Preprint, 2024. \href{https://arxiv.org/abs/2411.04872}{arXiv:2411.04872}.
Cited for research direction, not for a current model-performance claim.

\bibitem{cmms}
Filippo Calderoni, David Marker, Luca Motto Ros, and Assaf Shani.
\emph{Anti-classification results for groups acting freely on the line}.
\href{https://arxiv.org/abs/2010.08049}{arXiv:2010.08049}, version 2,
13 January 2023. See Lemma 2.2 and Section 7, especially Theorem 7.1,
Proposition 7.2, and Corollary 7.3.

\bibitem{marcone}
Alberto Marcone and Christian Rosendal.
The complexity of continuous embeddability between dendrites.
\emph{Journal of Symbolic Logic} \textbf{69}(3) (2004), 663--673.
The colored-order completeness theorem used here is also stated in
\cite[Theorem 7.1]{cmms}.

\bibitem{rast}
Richard Rast and Davender Singh Sahota.
The Borel complexity of isomorphism for o-minimal theories.
\emph{Journal of Symbolic Logic} \textbf{82}(2) (2017), 453--473.
\href{https://arxiv.org/abs/1408.5876}{arXiv:1408.5876}.

\bibitem{friedman}
Harvey Friedman and Lee Stanley.
A Borel reducibility theory for classes of countable structures.
\emph{Journal of Symbolic Logic} \textbf{54}(3) (1989), 894--914.

\bibitem{laver}
Richard Laver.
On Fra\"iss\'e's order type conjecture.
\emph{Annals of Mathematics} \textbf{93}(1) (1971), 89--111.

\bibitem{miller}
Benjamin D. Miller.
\emph{An introduction to classical descriptive set theory}.
Lecture notes, available at
\url{https://glimmeffros.github.io/seminars/descriptive.pdf}.
See Proposition 1.4.28 for completeness of ill-founded trees.

\bibitem{gonshor}
Harry Gonshor.
\emph{An Introduction to the Theory of Surreal Numbers}.
London Mathematical Society Lecture Note Series 110,
Cambridge University Press, 1986.

\bibitem{vde}
Lou van den Dries and Philip Ehrlich.
Fields of surreal numbers and exponentiation.
\emph{Fundamenta Mathematicae} \textbf{167}(2) (2001), 173--188;
erratum, \textbf{168} (2001), 295--297.
\href{https://doi.org/10.4064/fm167-2-3}{doi:10.4064/fm167-2-3}.

\bibitem{marker}
David Marker.
\emph{Model Theory: An Introduction}.
Graduate Texts in Mathematics 217, Springer, 2002.
Used for standard real-closed-field quantifier elimination and model
completeness.

\bibitem{proveit}
Vladimir Reshetnikov.
\emph{ProveIt}, surreal-number development.
Repository snapshot
\texttt{7c0f2d9f92c3d51ec85703bed1022924a4b7359b}, inspected
3 October 2026.
\url{https://github.com/VladimirReshetnikov/ProveIt}.
Relevant directory: \texttt{Algebra/SurrealNumbers/}.
Repository status statements are not independent verification of this
manuscript.

\bibitem{sse}
\emph{Self-Embeddings of the Surreal Numbers}.
Merged AI-assisted research report in ProveIt, at the same snapshot as
\cite{proveit}. Directory
\texttt{Algebra/SurrealNumbers/docs/surreal/}\allowbreak
\texttt{surreal-self-embeddings/}.
Its README explicitly distinguishes the report from formalized results.

\end{thebibliography}
\end{document}
